\newpage
\columnratio{0.6, 0.4}
\begin{paracol}{2}
\insimgcol{HRCP/005_folio_02_verso_1_B_SUPINE_Seq.jpg}

\asananamestmp{vṛṣapādakṣepāsanam}{Pawing the Leg like a Bull Pose}{ವೃಷಪಾದಕ್ಷೇಪಾಸನ}
\switchcolumn
\vfill
\begin{sans}
uttānaṃ śayanaṃ kṛtvā aṅgulibhiḥ kandharāṃ baddhvā kūrparau militvā\break nitambena bhūmiṃ spṛṣṭvā ekaṃ pādaṃ prasārya ekaikena pādena savyadakṣiṇaṃ bhrāmayitvā vṛṣapādakṣepaṃ bhavati\!॥1\!॥
\end{sans}

%1
%\incanum
\asana{[The yogi] should lie supine,  bind the neck with the fingers, join the elbows, touch the buttocks on the ground, extend one leg and rotate the other leg to the left and right. [This] is the \engname{Pawing the Leg like a Bull [Pose]} (\engiast{vṛṣapādakṣepa}).}
{ಯೋಗಿಯು ಮಲಗಿದ ಸ್ಥಿತಿಯಲ್ಲಿ ತನ್ನ ಬೆರಳುಗಳಿಂದ ಕುತ್ತಿಗೆಯನ್ನು ಹಿಡಿದು, ಮೊಣಕೈಗಳನ್ನು ಜೋಡಿಸಿರಬೇಕು, ಪೃಷ್ಠದ ಭಾಗವು ನೆಲಕ್ಕೆ ತಗಲುವಂತಿರಬೇಕು. ಒಂದು ಕಾಲನ್ನು ಚಾಚಿ, ಮತ್ತೊಂದು ಕಾಲನ್ನು ಎಡದಿಂದ ಬಲಕ್ಕೆ ತಿರುಗಿಸಬೇಕು. ಇದನ್ನು \tbf{ವೃಷಭ ತನ್ನ ಕಾಲನ್ನು ಕೆರೆಯುವ ಭಂಗಿ} ಎನ್ನಲಾಗುತ್ತದೆ (\tit{ವೃಷಪಾದ\-ಕ್ಷೇ\-ಪಾಸನ}).}
\vfill
\end{paracol}
\vskip 0.8cm

\columnratio{0.4, 0.6}

\begin{paracol}{2}
\vfill
\begin{sans}
uttānaṃ śayanaṃ kṛtvā pādau militvā prasārya nitambaṃ bhūmau spṛṣṭvā hastābhyāṃ kandharāṃ\endnote{kandharāṃ] P, kandharaṃ M} baddhvā kumbhakaṃ kṛtvā tiṣṭhet\!। parighāsanaṃ bhavati\!॥2\!॥
\end{sans}

%2.
%\incanum
\asana{[The yogi] should lie supine, join and extend the legs, touch the buttocks on the ground, bind the neck with the hands, hold the breath and remain thus. [This] is the \engname{Iron-bar Pose} (\engiast{parighāsana}).}
{ಮಲಗಿದ ಸ್ಥಿತಿಯಲ್ಲಿ ಕಾಲುಗಳನ್ನು ಜೋಡಿಸಿ ಚಾಚಬೇಕು. ಪೃಷ್ಠವು ನೆಲಕ್ಕೆ ತಗಲಿರಲಿ. ಕೈಗಳಿಂದ ಕುತ್ತಿಗೆಯನ್ನು ಹಿಡಿದು, ಉಸಿರನ್ನು ಬಿಗಿಹಿಡಿದು ಹಾಗೆಯೇ ಇರಬೇಕು. ಇದಕ್ಕೆ \tbf{ಕಬ್ಬಿಣದ ಕಂಬಿಯಂತಹ ಭಂಗಿ} ಎನ್ನಲಾಗುತ್ತದೆ (\tit{ಪರಿಘಾಸನ}).}
\vfill
\switchcolumn

\insimgcol{HRCP/006_folio_02_verso_2.jpg}

\asananamestmp{parighāsanaṃ}{Iron-bar Pose}{ಪರಿಘಾಸನ}

\end{paracol}


\newpage

\columnratio{0.6, 0.4}
\begin{paracol}{2}
\insimgcol{HRCP/007_folio_03_recto_3.jpg}

\asananamestmp{paraśvadhāsanaṃ}{Hatchet Pose}{ಪರಶ್ವಾಸನ}
\switchcolumn

\begin{sans}
uttānaṃ śayanaṃ kūrparadvayaṃ nābhau sthāpayitvā ekaikaṃ hastaṃ prasārya nāsikāyāmaṅguṣṭhapradeśena dhṛtvā tallakṣyeṇa jaghana\-pradeśena dhṛtvā sthāpayet\!। paraśvāsanaṃ\endnote{paraśvāsanaṃ] M(ac) P, paraśvadhāsanaṃ M(pc)} bhavati\!॥3\!॥
\end{sans}

%3.
\asana{[The yogi] should lie supine, place the elbows on the navel, extend one hand at a time, hold the nose by the thumb [of the other hand] with the gaze on it, while supporting [the position] with the region of the hips, and remain thus. [This] is the \engname{Hatchet Pose} (\engiast{paraśvāsana}).}
{ಮಲಗಿದ ಸ್ಥಿತಿಯಲ್ಲಿ ಮೊಣಕೈಗಳನ್ನು ನಾಭಿಯ (ಹೊಕ್ಕುಳ) ಮೇಲೆ ಇರಿಸಬೇಕು. ಒಂದೊಂದೇ ಕೈಯನ್ನು ಚಾಚುತ್ತಾ, ಒಂದು ಕೈಯ ಹೆಬ್ಬೆರಳಿನಿಂದ ಮೂಗಿನ ತುದಿಯನ್ನು ಹಿಡಿದು ಅದರತ್ತಲೇ ನೋಡಬೇಕು. ಇಡೀ ಶರೀರವು ಸೊಂಟದ ಭಾಗದ ಮೇಲೆ ಸಮತೋಲನದಲ್ಲಿರಬೇಕು. ಇದನ್ನು \tbf{ಕೊಡಲಿಯ ಭಂಗಿ} ಎನ್ನಲಾಗುತ್ತದೆ (\tit{ಪರಶ್ವಾಸನ}).}

\end{paracol}
\vskip 0.8cm

\columnratio{0.45, 0.55}
\begin{paracol}{2}
\vfill
\begin{sans}
uttānaṃ śayanamekaikaṃ pādaṃ grīvāyāṃ vinyasya itarahastena pādāgraṃ gṛhītvā itarapādahastau lambīkṛtya tiṣṭhet\!। anantāsanaṃ bhavati\!॥4\!॥
\end{sans}

%4.
\asana{[The yogi] should lie supine, fix one foot on [the back of] the neck, grasp the toes of [that] foot with the other hand, and stretch out the other foot and hand. He should remain thus [and then repeat it with the other foot on the neck]. [This] is \engname{Ananta’s Pose} (\engiast{anantāsana}).}
{ಮಲಗಿದ ಸ್ಥಿತಿಯಲ್ಲಿ ಎಡ ಪಾದವನ್ನು ಕುತ್ತಿಗೆಯ ಹಿಂಭಾಗದಲ್ಲಿ ಇರಿಸಿ, ಅದೇ ಪಾದದ ಬೆರಳುಗಳನ್ನು ಬಲಗೈಯಿಂದ ಹಿಡಿಯಬೇಕು. ಉಳಿದ ಒಂದು ಕಾಲು ಮತ್ತು ಕೈಯನ್ನು ನೇರವಾಗಿ ಚಾಚಬೇಕು. ಇದೇ ರೀತಿ ಮತ್ತೊಂದು ಕಾಲಿನಿಂದಲೂ ಮಾಡಬೇಕು. ಇದು \tbf{ಅನಂತಾಸನ}.}
\vfill
\switchcolumn

\insimgcol{HRCP/008_folio_03_recto_4.jpg}

\asananamestmp{anantāsanaṃ}{Ananta’s Pose}{ಅನಂತಾಸನ}

\end{paracol}

\newpage

\columnratio{0.55, 0.45}
\begin{paracol}{2}
\insimgcol{HRCP/009_folio_03_verso_5.jpg}

\asananamestmp{aṅkuśāsanaṃ}{Goad Pose}{ಅಂಕುಶಾಸನ}
\switchcolumn

\begin{sans}
uttānaṃ śayīta\endnote{śayīta~] \textit{emend.}, śayātaṃ M.}\!। ekaikaṃ pādaṃ grīvāyāṃ kurvannitarahastaṃ karṇamūle sthāpayet\!। tasyaiva hastasya kūrparaṃ bhūmau nidhāya itarahastapādau  saralīkṛtya tiṣṭhet\!। aṅkuśāsanaṃ bhavati\!॥5\!॥
\end{sans}

%5.
\asana{[The yogi] should lie supine. Putting one foot on the [back of the] neck, he should place the other hand on the root of the ear. Having placed the elbow of this arm on the ground, while making the other hand and leg straight, he should remain thus [and then repeat it with the other foot on the neck]. [This] is the \engname{Goad Pose} (\engiast{aṅkuśāsana}).}
{ಮಲಗಿದ ಸ್ಥಿತಿಯಲ್ಲಿ ಎಡ ಪಾದವನ್ನು ಕುತ್ತಿಗೆಯ ಹಿಂದೆ ಇರಿಸಿ, ಬಲಗೈಯನ್ನು ಕಿವಿಯ ಬುಡದ ಹತ್ತಿರ ಇಡಬೇಕು. ಆ ತೋಳಿನ ಮೊಣಕೈ ನೆಲಕ್ಕೆ ತಗಲಿರಲಿ. ಉಳಿದ ಒಂದು ಕೈ ಮತ್ತು ಕಾಲು ನೇರವಾಗಿರಬೇಕು. ಇದೇ ರೀತಿ ಎರಡೂ ಕಡೆ ಅಭ್ಯಾಸ ಮಾಡಬೇಕು. ಇದನ್ನು \tbf{ಅಂಕುಶಾಸನ} ಎನ್ನಲಾಗುತ್ತದೆ.}
\end{paracol}
\vskip 0.8cm

\columnratio{0.45, 0.55}
\begin{paracol}{2}
\vfill
\begin{sans}
śavavaccharīraṃ\endnote{śavavac°~] P, śavava M.} saṃsthāpya jānunī saṃmīlya nābhāvānīya hastābhyāṃ kandharāṃ baddhvā bhrāmayet\!। śvottānaṃ bhavati\!॥6\!॥
\end{sans}

%6.
\asana{Having placed the body [supine] like a corpse, [the yogi] should join the knees together, bring [them] to the navel, clasp the neck with the hands and rotate [the legs]. [This] is the \engname{Upturned Dog [Pose]} (\engiast{śvottāna}).}
{ಶವದಂತೆ ನೇರವಾಗಿ ಮಲಗಿ, ಎರಡು ಮೊಣಕಾಲುಗಳನ್ನು ಜೋಡಿಸಿ ನಾಭಿಯ ಹತ್ತಿರಕ್ಕೆ ತರಬೇಕು. ಕೈಗಳಿಂದ ಕುತ್ತಿಗೆಯನ್ನು ತಬ್ಬಿಕೊಂಡು ಕಾಲುಗಳನ್ನು ತಿರುಗಿಸಬೇಕು. ಇದನ್ನು \tbf{ಮೇಲ್ಮುಖವಾದ ನಾಯಿಯ ಭಂಗಿ} ಎಂದು ಕರೆಯಲಾಗುತ್ತದೆ (\tit{ಶ್ವೋತ್ತಾನಾಸನ}).}
\vfill
\switchcolumn

\insimgcol{HRCP/010_folio_03_verso_6.jpg}

\asananamestmp{śvottānaṃ}{Upturned Dog [Pose]}{ಶ್ವೋತ್ತಾನಾಸನ}
\end{paracol}


\newpage

\columnratio{0.55, 0.45}
\begin{paracol}{2}
\insimgcol{HRCP/011_folio_04_recto_7.jpg}

\asananamestmp{mārjārottānaṃ}{Upturned Cat [Pose]}{ಮಾರ್ಜಾರೊತ್ತಾನಾಸನ}
\switchcolumn

\begin{sans}
śvottānavatsaṃsthitiṃ kṛtvā jānudvayaṃ paryāyeṇa karṇayoḥ saṃspṛśet\!। mārjārottānaṃ bhavati\!॥7\!॥
\end{sans}

%7.
\asana{Having positioned [himself] as in the upturned dog [pose, the yogi] should touch the knees with his ears in turn. [This is] the \engname{Upturned Cat [Pose]} (\engiast{mārjārottāna}).}
{ಶ್ವೋತ್ತಾನಾಸನದಲ್ಲಿರುವಂತೆಯೇ ಇದ್ದು, ಮೊಣಕಾಲುಗಳನ್ನು ಸರದಿಯಂತೆ ಕಿವಿಗಳಿಗೆ ತಾಗಿಸಬೇಕು. ಇದನ್ನು  \tbf{ಮೇಲ್ಮುಖವಾದ ಬೆಕ್ಕಿನ ಭಂಗಿ} ಎಂದು ಕರೆಯುತ್ತಾರೆ (\tit{ಮಾರ್ಜಾರೋತ್ತಾನಾಸನ}).}

\end{paracol}
\vskip 0.5cm

\columnratio{0.45, 0.55}
\begin{paracol}{2}
\vfill
\begin{sans}
atha uttānaṃ śayanaṃ kṛtvā pādatalābhyāṃ bhūmiṃ dhṛtvā uttiṣṭhet\!। vṛkāsanaṃ\!॥8\!॥
\end{sans}

%8.
\asana{Then, having lain supine, [the yogi] should grip  the ground with the soles of the feet and  stand up. [This] is the \engname{Wolf Pose} (\engiast{vṛkāsana}).}
{ಮಲಗಿದ ಸ್ಥಿತಿಯಲ್ಲಿ ಪಾದಗಳನ್ನು ನೆಲಕ್ಕೆ ಗಟ್ಟಿಯಾಗಿ ಒತ್ತಿ ಹಿಡಿದು ನೇರವಾಗಿ ಎದ್ದು ನಿಲ್ಲಬೇಕು. ಇದನ್ನು \tbf{ತೋಳದ ಭಂಗಿ} ಎನ್ನಲಾಗುತ್ತದೆ (\tit{ವೃಕಾಸನ}).}
\vfill
\switchcolumn

\insimgcol{HRCP/012_folio_04_recto_8.jpg}

\asananamestmp{vṛkāsanaṃ}{Wolf Pose}{ವೃಕಾಸನ}
\end{paracol}

\newpage

\columnratio{0.55, 0.45}
\begin{paracol}{2}
\insimgcol{HRCP/013_folio_04_verso_9.jpg}

\asananamestmp{trikūṭaṃ}{Mountain with Three Peaks [Pose]}{ತ್ರಿಕೂಟ ಪರ್ವತಾಸನ}
\switchcolumn

\begin{sans}
uttānaśayanaṃ pādatalābhyāṃ bhūmiṃ dhṛtvā kūrparau avanau nidhāya pṛṣṭhapradeśamūrdhvamunnayet\!। trikūṭaṃ bhavati\!॥9\!॥
\end{sans}

%9.
\asana{Lying supine, [the yogi] should grip the ground with the soles of the feet,  put the elbows on the ground and lift up his back. [This] is the \engname{Mountain with Three Peaks [Pose]} (\engiast{trikūṭa}).}
{ಮಲಗಿದ ಸ್ಥಿತಿಯಲ್ಲಿ ಪಾದಗಳನ್ನು ನೆಲಕ್ಕೆ ಊರಿ, ಮೊಣಕೈಗಳನ್ನು ನೆಲದ ಮೇಲೊತ್ತಿ ಬೆನ್ನಿನ ಭಾಗವನ್ನು ಮೇಲೆತ್ತಬೇಕು. ಇದನ್ನು \tbf{ತ್ರಿಕೂಟ ಪರ್ವತಾಸನ} ಎನ್ನಲಾಗುತ್ತದೆ.}
\end{paracol}
\vskip 0.2cm

\long\def\asana#1#2{%
\refstepcounter{anum}
\begin{eng}
\begin{enumerate}[leftmargin=*,labelindent=0pt,widest=99, nosep]
\item[\asananum]
#1
\end{enumerate}
\end{eng}
\begin{kan}
\begin{enumerate}[leftmargin=*,labelindent=0pt,widest=99, nosep]
\item[\asananum]
#2
\end{enumerate}
\end{kan}
}
\columnratio{0.45, 0.55}
\begin{paracol}{2}
\vfill
\begin{sans}
uttānaśayanaṃ pādatale uttānahastayoḥ sthāpayet\!। pṛṣṭhabhāgamunnamayet\!। kamaṭhapīṭhaṃ bhavati\!॥10\!॥
\end{sans}
%\bigskip

%10.
\asana{Lying supine, [the yogi] should place the soles of the feet on the upturned hands [and] raise the back part [of the body from the ground]. [This] is the \engname{Tortoise’s Seat} (\engiast{kamaṭhapīṭha}).\endnote{This pose is called the Monkey’s Pose' (\textit{markaṭapīṭha}) in the Pune manuscript.}}
{ಮಲಗಿದ ಸ್ಥಿತಿಯಲ್ಲಿ ಅಂಗೈಗಳನ್ನು ಮೇಲ್ಮುಖವಾಗಿ ಇಟ್ಟು ಅವುಗಳ ಮೇಲೆ ಪಾದಗಳನ್ನು ಇರಿಸಬೇಕು, ನಂತರ ಶರೀರದ ಹಿಂಭಾಗವನ್ನು ನೆಲದಿಂದ ಮೇಲೆತ್ತಬೇಕು. ಇದು  \tbf{ಆಮೆಯ ಭಂಗಿ} (\tit{ಕಮಠಪೀಠಾಸನ}).}
\vfill
\switchcolumn

\insimgcol{HRCP/014_folio_04_verso_10.jpg}

\asananamestmp{kamaṭhapīṭhaṃ}{Tortoise’s Seat}{ಕಮಠಪೀಠಾಸನ}
\end{paracol}


\newpage

\columnratio{0.55, 0.45}
\begin{paracol}{2}
\insimgcol{HRCP/015_folio_05_recto_11.jpg}

\asananamestmp{naukā}{Boat [Pose]}{ನೌಕಾಸನ}
\switchcolumn
\begin{sans}
uttānaśayanaṃ kūrparābhyāṃ dharaṇīmavaṣṭabhya hastau nitambe nidhāya śiraḥūrujaṅghāpādāndaṇḍavaddhārayet\!। naukā bhavati\!॥11\!॥
\end{sans}
%\bigskip

%11.
\asana{Lying supine, [the yogi] should support [himself] with the elbows on the ground, place the hands on the buttocks and hold his head, thighs, lower legs and feet [straight] like a stick. [This] is the \engname{Boat [Pose]} (\engiast{naukā}).}
{ಮಲಗಿದ ಸ್ಥಿತಿಯಲ್ಲಿ ಮೊಣಕೈಗಳನ್ನು ನೆಲಕ್ಕೆ ಊರಿ ಆಧಾರ ಪಡೆದು, ಕೈಗಳನ್ನು ಪೃಷ್ಠದ ಮೇಲಿರಿಸಬೇಕು. ತಲೆ, ತೊಡೆ ಮತ್ತು ಕಾಲುಗಳನ್ನು ದಂಡದಂತೆ (ಕೋಲಿನಂತೆ) ನೇರವಾಗಿ ಹಿಡಿಯಬೇಕು. ಇದನ್ನು  \tbf{ದೋಣಿಯ ಭಂಗಿ} ಎನ್ನಲಾಗುತ್ತದೆ (\tit{ನೌಕಾಸನ}).}
\end{paracol}
\vskip 0.8cm

\columnratio{0.45, 0.55}
\begin{paracol}{2}
\vfill
\begin{sans}
naukāvatsthitvā ūrdhvaṃ pādāgradvayaṃ nayet\!। tiryaṅnaukāsanaṃ bhavati\!॥12\!॥
\end{sans}
%\bigskip

%12.
\asana{While in the Boat [Pose, the yogi] should bring the toes of the feet upwards. This is the \engname{Diagonal Boat Pose} (\engiast{tiryaṅnaukāsana}).}
{ನೌಕಾಸನದಲ್ಲಿರುವಾಗಲೇ ಕಾಲ್ಬೆರಳುಗಳನ್ನು ಮೇಲ್ಮುಖವಾಗಿ ತರಬೇಕು. ಇದನ್ನು \tbf{ಓರೆಯಾದ ದೋಣಿ ಭಂಗಿ} ಎನ್ನಲಾಗುತ್ತದೆ (\tit{ತಿರ್ಯನ್ನೌಕಾಸನ}).}
\vfill
\switchcolumn

\insimgcol{HRCP/016_folio_05_recto_12.jpg}

\asananamestmp{tiryaṅnaukāsanaṃ}{Diagonal Boat Pose}{ತಿರ್ಯನ್ನೌಕಾಸನ}
\end{paracol}


\newpage

\insimg[0.75]{HRCP/017_folio_05_verso_13.jpg}

\asananamestmpscale{dhvajāsanaṃ}{Banner Pose}{ಧ್ವಜಾಸನ}{0.75}
\vspace{0.5cm}

\begin{adjustwidth}{3cm}{3cm}
\begin{sans}
tiryaṅnaukāvadgrīvāpṛṣṭhakūrparaiḥ bhūmimavaṣṭabhya mastakalakṣeṇa  pādāgradvayamūrdhvaṃ nayet\!। dhvajāsanaṃ bhavati\!॥13\!॥
\end{sans}
\begin{multicols}{2}
%13.
\asana{[Staying] as though in the Diagonal Boat [Pose] and supporting [himself] with the neck, back and elbows on the ground, [the yogi] should take up the toes of the feet, pointing [them] towards the head. [This] is the \engname{Banner Pose} (\engiast{dhvajāsana}).}
{ತಿರ್ಯನ್ನೌಕಾಸನದಲ್ಲಿರುವಂತೆಯೇ ಇದ್ದು ಕುತ್ತಿಗೆ, ಬೆನ್ನು ಮತ್ತು ಮೊಣಕೈಗಳನ್ನು ನೆಲಕ್ಕೆ ಊರಬೇಕು. ಕಾಲ್ಬೆರಳುಗಳನ್ನು ತಲೆಯ ಕಡೆಗೆ ಬರುವಂತೆ ಹಿಡಿಯಬೇಕು. ಇದು \tbf{ಬಾವುಟದ ಭಂಗಿ} (\tit{ಧ್ವಜಾಸನ}).}
\end{multicols}
\end{adjustwidth}

\newpage

\insimg[0.75]{HRCP/018_folio_05_verso_14.jpg}

\asananamestmpscale{narakāsanaṃ}{Hell Pose}{ನರಕಾಸನ}{0.75}
\vspace{0.5cm}

\begin{adjustwidth}{3cm}{3cm}
\begin{sans}
grīvākaṇṭhena\endnote{°kaṇṭhena~] P, °kaṃṭhakena M.} bhūmiṃ viṣṭabhya pādāgradvayamūrdhvamunnayet\!। narakāsanaṃ bhavati\!॥14\!॥
\end{sans}

\begin{multicols}{2}
%14.
\asana{Having [firmly]  planted the nape of the neck on the ground, [the yogi] should raise up the toes. [This] is the \engname{Hell Pose} (\engiast{narakāsana}).}
{ಕುತ್ತಿಗೆಯ ಹಿಂಭಾಗವನ್ನು ನೆಲಕ್ಕೆ ಗಟ್ಟಿಯಾಗಿ ಒತ್ತಿ ಹಿಡಿದು, ಕಾಲುಗಳನ್ನು ಮೇಲೆತ್ತಿ ಕಾಲ್ಬೆರಳುಗಳನ್ನು ತಲೆಯ ಕಡೆಗೆ ಬರುವಂತೆ ಹಿಡಿಯಬೇಕು. ಇದನ್ನು \tbf{ನರಕಾಸನ} ಎನ್ನಲಾಗುತ್ತದೆ.}
\end{multicols}
\end{adjustwidth}

\newpage

\columnratio{0.53, 0.47}
\begin{paracol}{2}
\insimgcol{HRCP/018_folio_06_recto_15.jpg}

\asananamestmp{lāṅgalāsanaṃ}{Plough Pose}{ಲಾಂಗಲಾಸನ}

\switchcolumn

\begin{sans}
narakāsane sthitvā nāsikāpradeśe bhūmau pādapṛṣṭhe sthāpya hastadvayaṃ saṃmīlya lambīkuryādgrīvāpradeśena bhūmiṃ karṣayet\!। lāṅgalāsanaṃ bhavati\!॥15\!॥
\end{sans}

%15.
\asana{While in the Hell Pose, [the yogi] should place the upper side of the feet on the ground in the region of the nose, join and extend the hands, and plough the ground with the neck region. [This] is the \engname{Plough Pose} (\engiast{lāṅgalāsana}).}
{ನರಕಾಸನದಲ್ಲಿರುವಾಗಲೇ ಪಾದಗಳ ಮೇಲ್ಭಾಗವನ್ನು ಮೂಗಿನ ಹತ್ತಿರ ನೆಲಕ್ಕೆ ತಾಗಿಸಬೇಕು. ಕೈಗಳನ್ನು ಜೋಡಿಸಿ ಚಾಚುತ್ತಾ, ಕುತ್ತಿಗೆಯ ಭಾಗದಿಂದ ನೆಲವನ್ನು ನೇಗಿಲಿನಿಂದ ಉಳುಮೆ ಮಾಡುವಂತೆ ಒತ್ತಬೇಕು. ಇದು  \tbf{ನೇಗಿಲಿನ ಭಂಗಿ} (\tit{ಲಾಂಗಲಾಸನ}).}
\end{paracol}
\vskip 0.3cm

\columnratio{0.45, 0.55}
\begin{paracol}{2}
\vfill
\begin{sans}
uttānaśayanaṃ hastatalābhyāṃ bhūmimavaṣṭabhya pādatalābhyāṃ bhūmiṃ dhṛtvā nābhipradeśamūrdhvaṃ kuryāt\!। paryaṅkāsanaṃ bhavati\!॥16\!॥
\end{sans}

%16. 
\asana{Lying supine, [the yogi] should support himself with the palms on the ground, grip the ground with the soles of the feet, [and] raise up the region of the navel. [This] is the \engname{Sofa Pose} (\engiast{paryaṅkāsana}).}
{ಮಲಗಿದ ಸ್ಥಿತಿಯಲ್ಲಿ ಅಂಗೈ ಮತ್ತು ಪಾದಗಳನ್ನು ನೆಲಕ್ಕೆ ಊರಿ, ಹೊಕ್ಕುಳಿನ (ನಾಭಿ) ಭಾಗವನ್ನು ಮೇಲೆತ್ತಬೇಕು. ಇದನ್ನು \tbf{ಮಂಚದ ಭಂಗಿ} ಎನ್ನಲಾಗುತ್ತದೆ (\tit{ಪರ್ಯಂಕಾಸನ}).}
\vfill
\switchcolumn
\insimgcol{HRCP/019_folio_06_recto_16.jpg}

\asananamestmp{paryaṅkāsanaṃ}{Sofa Pose}{ಪರ್ಯಂಕಾಸನ}
\end{paracol}

\newpage


\columnratio{0.5, 0.5}
\begin{paracol}{2}
~\vskip 3.3cm
\begin{sans}
paryaṅkāsane sthitvā hastapādau saṃmīlayet\endnote{saṃmīlayet~] P, saṃmīlya M.}\!। vetrāsanaṃ bhavati\!॥17\!॥
\end{sans}

%17.
\asana{While in the Sofa Pose, [the yogi] should join his hands and feet. [This] is the \engname{Cane Pose} (\engiast{vetrāsana}).}
{ಪರ್ಯಂಕಾಸನದಲ್ಲಿರುವಾಗಲೇ ಕೈ ಮತ್ತು ಪಾದಗಳನ್ನು ಪರಸ್ಪರ ಜೋಡಿಸಬೇಕು. ಇದನ್ನು \tbf{ಬಿದಿರಿನ ಭಂಗಿ} ಎನ್ನಲಾಗುತ್ತದೆ (\tit{ವೇತ್ರಾಸನ}).}
\bigskip

\insimgcol{HRCP/020_folio_06_verso_17.jpg}

\asananamestmp{vetrāsanaṃ}{Cane Pose}{ವೇತ್ರಾಸನ}


\switchcolumn

\insimgcol{HRCP/021_folio_06_verso_18.jpg}

\asananamestmp{kandukāsanaṃ}{Ball Pose}{ಕಂದೂಕಾಸನ}
\bigskip

\begin{sans}
vetrāsane sthitvā hastapādānniṣkṛṣya\endnote{niṣkṛṣya~] \textit{emend.}, niṣkṛṣyam M.} ūrdhvaṃ nayetpṛṣṭhavaṃśena bhūmiṃ pīḍayet\endnote{pīḍayet~] \textit{conj.}, pothayet M.}\!। kandukāsanaṃ bhavati\!॥18\!॥
\end{sans}

%18. 
\asana{While in the Cane Pose, [the yogi] should release his hands and feet, take [them] upwards [and] press the ground with his backbone. [This] is the \engname{Ball Pose} (\engiast{kandukāsana}).}
{ವೇತ್ರಾಸನದಲ್ಲಿರುವಾಗಲೇ ಕೈ ಮತ್ತು ಪಾದಗಳನ್ನು ಮುಕ್ತಗಳಿಸಿ, ಮೊಣಕಾಲುಗಳನ್ನು ಜೋಡಿಸಿ, ಕೈಗಳಿಂದ ತಬ್ಬಿ, ಬೆನ್ನೆಲುಬಿನಿಂದ ನೆಲವನ್ನು ಒತ್ತಬೇಕು. ಇದನ್ನು  \tbf{ಚೆಂಡಿನ ಭಂಗಿ} ಎನ್ನಲಾಗುತ್ತದೆ (\tit{ಕಂದೂಕಾಸನ}).}

\end{paracol}

\newpage

\columnratio{0.5, 0.5}
\begin{paracol}{2}
\insimgcol{HRCP/022_folio_07_recto_19.jpg}

\asananamestmp{uttānakūrmaṃ}{Upturned Turtle [Pose]}{ಉತ್ತಾನಕೂರ್ಮಾಸನ}

\switchcolumn

\begin{sans}
ekasminnūruṇi ekaṃ pādaṃ kṛtvā anyasminnūruṇi anyaṃ pādaṃ kṛtvā padmāsanaṃ bhavati\!। ūrujānvorantarayoḥ\endnote{antarayoḥ~] P, aṃtayoḥ M.} hastau praveśya\endnote{praveśya~]  \textit{emend.}, praviśya M.} kandharāṃ baddhvā uttānaṃ tiṣṭhet\!। uttānakūrmaṃ bhavati\!॥19\!॥
\end{sans}

%19.
\asana{Placing one foot on one thigh and the other foot on the other thigh is the Lotus Pose. Having passed the hands in the gaps between the thighs and knees, [the yogi] should clasp the neck  and lie on his back. [This] is the \engname{Upturned Turtle [Pose]} (\engiast{uttānakūrma}).}{ಒಂದು ಪಾದವನ್ನು ಇನ್ನೊಂದು ತೊಡೆಯ ಮೇಲಿಟ್ಟು ಪದ್ಮಾಸನ ಹಾಕಬೇಕು. ನಂತರ ಮೊಣಕಾಲು ಮತ್ತು ತೊಡೆಗಳ ಮಧ್ಯದ ಸಂದಿನಿಂದ ಕೈಗಳನ್ನು ತೂರಿಸಿ ಕುತ್ತಿಗೆಯನ್ನು ಹಿಡಿದು ಬೆನ್ನ ಮೇಲೆ ಮಲಗಬೇಕು. ಇದು \tbf{ಮೇಲ್ಮುಖವಾದ ಆಮೆ ಭಂಗಿ} (\tit{ಉತ್ತಾನಕೂರ್ಮಾಸನ}).}

\end{paracol}
\vskip 0.8cm

\columnratio{0.5, 0.5}
\begin{paracol}{2}
\vfill
\begin{sans}
jaṅghāpṛṣṭhe\endnote{°pṛṣṭhe~] \textit{emend.}, °pṛṣṭhau M.} bhūmau nidhāya jaṅghodarayoḥ ūruṇī saṃsthāpya pṛṣṭhavaṃśaṃ vāraṃ vāraṃ spṛśyet\endnote{spṛśyet~] \textit{corr.},  spṛśet M.}\!। viratāsanaṃ\endnote{The manuscript has the letter \textit{lā} written above the \textit{tā} of \textit{viratāsanaṃ}. The scribe has not indicated whether this was intended as a correction. However, the name \textit{viralāsana} (`the rare pose') is rather odd.}  bhavati\!॥20\!॥
\end{sans}

%20.
\asana{Placing the shins and back on the ground and positioning the thighs on the calves, [the yogi] should touch his backbone [on the ground] again and again. [This] is the \engname{Pose for One who has Ceased [from Worldly Activities]} (\engiast{viratāsana}).}{ಕಾಲುಗಳ ಕೆಳಭಾಗ ಮತ್ತು ಬೆನ್ನನ್ನು ನೆಲಕ್ಕೆ ಊರಿ, ತೊಡೆಗಳನ್ನು ಪಾದದ ಹಿಮ್ಮಡಿಯ ಮೇಲಿರಿಸಿ, ಬೆನ್ನೆಲುಬನ್ನು ಪದೇ ಪದೇ ನೆಲಕ್ಕೆ ತಾಗಿಸಬೇಕು. ಇದನ್ನು \tbf{ಲೌಕಿಕ ವ್ಯವಹಾರಗಳಿಂದ ದೂರಾದವನ ಭಂಗಿ} ಎನ್ನಲಾಗುತ್ತದೆ (\tit{ವಿರತಾಸನ}).}
\vfill
\switchcolumn

\insimgcol{HRCP/023_folio_07_recto_20.jpg}

\asananamestmp{viratāsanaṃ}{Pose for One who has Ceased [from Worldly Activities]}{ವಿರತಾಸನ}

\end{paracol}
\newpage

\columnratio{0.55, 0.45}
\begin{paracol}{2}
\insimgcol{HRCP/024_folio_07_verso_21.jpg}

\asananamestmp{dṛṣadāsanaṃ}{Grindstone Pose}{ದೃಶದಾಸನ}
\switchcolumn

\begin{sans}
uttānaśayanaṃ kṛtvā jānunī hṛdaye 'vaṣṭabhya jaṅghāsahita ūruṇī karadvayena\endnote{°dvayena~] \textit{conj.}, °dvaye M.}  baddhvā savyāpasavyaṃ loḍayet\!। dṛṣadāsanaṃ bhavati\!॥21\!॥
\end{sans}

%21.
\asana{Lying supine, [the yogi] should support the knees on the chest, bind the thighs and shins with the hands, and rock to the left and right. [This] is the \engname{Grindstone Pose} (\engiast{dṛṣadāsana}).}
{ಮಲಗಿದ ಸ್ಥಿತಿಯಲ್ಲಿ ಮೊಣಕಾಲುಗಳನ್ನು ಎದೆಗೆ ತಾಗಿಸಿ, ಕೈಗಳಿಂದ ತೊಡೆ ಮತ್ತು ಕಾಲುಗಳನ್ನು ಬಿಗಿಯಾಗಿ ಅಪ್ಪಿಕೊಳ್ಳಬೇಕು. ನಂತರ ಎಡಕ್ಕೆ ಮತ್ತು ಬಲಕ್ಕೆ ಉರುಳಾಡಬೇಕು. ಇದನ್ನು \tbf{ಗಿರಣಿ ಅಥವಾ ಬೀಸುವ ಕಲ್ಲಿನ ಭಂಗಿ} ಎನ್ನಲಾಗುತ್ತದೆ (\tit{ದೃಶದಾಸನ}).}
\end{paracol}
\vskip 0.6cm

\columnratio{0.45, 0.55}
\begin{paracol}{2}
\vfill
\begin{sans}
uttānaśayanaṃ pādau śiraḥ ullaṅghya bhūmau sthāpayitvā nyubjaṃ bhūyate\endnote{bhūyate~] \textit{emend.}, bhūyāt M.}\!। nyubjaṃ śayanaṃ pṛṣṭhaṃ bhūmau nidhāya punaḥ punaḥ paryāyeṇa kartavyam\!। luṇṭhanāsanaṃ bhavati\!॥22\!॥
\end{sans}

%22.
\asana{Lying supine, [the yogi] should pass the feet over the head, place them on the ground and turn onto his front. Lying prone, [the yogi] should place his back on the ground and turn over again and again. [This] is the \engname{Rolling Pose} (\engiast{luṇṭhanāsana}).\endnote{The spelling of the name of this \textit{āsana} in the Mysore manuscript (i.e., \textit{luṇṭhana}) is a variant of \textit{luṭhana} (the reading of P) attested in the Monier-Williams dictionary (s.v.\ 1899).}}{ಮಲಗಿದ ಸ್ಥಿತಿಯಲ್ಲಿ ಕಾಲುಗಳನ್ನು ತಲೆಯ ಮೇಲಿಂದ ತಂದು ನೆಲಕ್ಕೆ ತಾಗಿಸಿ, ನಂತರ ಮಗುಚಿ ಬಿದ್ದು ಹೊಟ್ಟೆಯ ಮೇಲೆ ಮಲಗಬೇಕು. ಮತ್ತೆ ಬೆನ್ನಿನ ಮೇಲೆ ಮಲಗುತ್ತಾ ಹೀಗೆ ಪದೇ ಪದೇ ಉರುಳಬೇಕು. ಇದನ್ನು \tbf{ಉರುಳುವ ಭಂಗಿ} ಎನ್ನಲಾಗುತ್ತದೆ (\tit{ಲುಂಠನಾಸನ}).}
\vfill
\switchcolumn
\insimgcol{HRCP/025_folio_07_verso_22_E_SUPINE_Seq.jpg}

\asananamestmp{luṇṭhanāsanaṃ}{Rolling Pose}{ಲುಂಠನಾಸನ}
\end{paracol}

\newpage

\insertFullImage{041_folio_11_verso_38.jpg}

\newpage

\pagesWithBorder

\null\vfill
{\leftskip=30pt
\begin{sans}
{\fontsize{18pt}{21.6pt}\selectfont \color{deepmaroon}{atha nyubjāsanāni\!॥}}
\end{sans}
\vskip 7pt

\begin{eng}
\prtheading[Prone Postures]{Now, the Prone Postures}
\end{eng}
\vskip 7pt

\begin{kan}
\prtheading[ಹೊಟ್ಟೆಯ ಮೇಲೆ ಮಗುಚಿ ಬಿದ್ದು ಮಾಡುವ ಆಸನಗಳು]{ಈಗ, ಹೊಟ್ಟೆಯ ಮೇಲೆ ಮಗುಚಿ ಬಿದ್ದು ಮಾಡುವ ಆಸನಗಳು}
\end{kan}\par}
\vfill

\newpage

\columnratio{0.55, 0.45}
\begin{paracol}{2}
\insimgcol{HRCP/026_folio_08_recto_23_B_PRONE.jpg}

\asananamestmp{saraṭāsanaṃ}{Lizard Pose}{ಸರಟಾಸನ}
\switchcolumn

\begin{sans}
nyubjaṃ śayanaṃ kṛtvā nābhiṃ bhūmau nidhāya stambhavaddhastābhyāṃ bhūmim avaṣṭabhya oṣṭhau nimīlya veṇuvatsūṃkṛtya tiṣṭhet\!। saraṭāsanaṃ bhavati\!॥23\!॥
\end{sans}

%23.
\asana{Lying prone, [the yogi] should place the navel on the ground and support [himself on] the ground with the forearms like pillars, join the lips, make the [sound] `\textit{sūṃ}' like a flute, and remain thus. [This] is the \engname{Lizard Pose} (\engiast{saraṭāsana}).}
{ಹೊಟ್ಟೆ ಕೆಳಗೆ ಮಾಡಿ ಮಲಗಿ, ಹೊಕ್ಕುಳನ್ನು ನೆಲಕ್ಕೆ ತಾಗಿಸಿ, ಮುಂಗೈಗಳನ್ನು ಕಂಬಗಳಂತೆ ನೆಲಕ್ಕೆ ಊರಿ ಆಧಾರ ಪಡೆಯಬೇಕು. ತುಟಿಗಳನ್ನು ಜೋಡಿಸಿ ಕೊಳಲಿನಂತೆ ‘ಸೂಂ’ ಎಂದು ಶಬ್ದ ಮಾಡುತ್ತಾ ಹಾಗೆಯೇ ಇರಬೇಕು. ಇದನ್ನು \tbf{ಓತಿಯ ಭಂಗಿ} ಎನ್ನಲಾಗುತ್ತದೆ (\tit{ಸರಟಾಸನ}).}
\end{paracol}
\vskip 0.8cm

\columnratio{0.45, 0.55}
\begin{paracol}{2}
\vfill
\begin{sans}
nyubjaṃ śayanaṃ kūrparau pārśvabhāgābhyāmūrdhvīkṛtya hastatalābhyāṃ bhūmimavaṣṭabhya vāraṃ vāramuḍyānaṃ kuryāt\endnote{kuryāt~] P, kṛtvā M.}\!। matsyāsanaṃ bhavati\!॥24\!॥
\end{sans}

%24.
\asana{Lying prone, [the yogi] should raise up the elbows by the sides [of the body,] and support [himself on] the ground with the palms of the hands. He should fly up again and again. [This] is the \engname{Fish Pose} (\engiast{matsyāsana}).}
{ಹೊಟ್ಟೆ ಕೆಳಗೆ ಮಾಡಿ ಮಲಗಿ, ಮೊಣಕೈಗಳನ್ನು ಶರೀರದ ಪಕ್ಕದಲ್ಲಿ ಮೇಲೆತ್ತಿ, ಅಂಗೈಗಳನ್ನು ನೆಲಕ್ಕೆ ಊರಬೇಕು. ನಂತರ ಪದೇ ಪದೇ ಮೇಲಕ್ಕೆ ಹಾರಲು ಯತ್ನಿಸಬೇಕು. ಇದು \tbf{ಮೀನಿನ ಭಂಗಿ} (\tit{ಮತ್ಸ್ಯಾಸನ}).}
\vfill
\switchcolumn

\insimgcol{HRCP/027_folio_08_recto_24.jpg}

\asananamestmp{matsyāsanaṃ}{Fish Pose}{ಮತ್ಸ್ಯಾಸನ}
\end{paracol}

\newpage

\columnratio{0.55, 0.45}
\begin{paracol}{2}
\insimgcol{HRCP/028_folio_08_verso_25.jpg}

\asananamestmp{gajāsanaṃ}{Elephant Pose}{ಗಜಾಸನ}
\switchcolumn

\begin{sans}
nyubjaṃ pādāgre bhūmau kṛtvā lambībhūya mastakāgre hastatalau nidhāya nitambamūrdhvamunnamayya nābhiṃ\endnote{nābhiṃ~] \textit{corr.}~: nābhi M.} lakṣya bhūmau nāsikāmānīya hastatalaparyantaṃ nayet\!। itthaṃ punaḥ punaḥ\endnote{punaḥ punaḥ~] P, punaḥ M.} kuryāt\!। gajāsanaṃ bhavati\!॥25\!॥
\end{sans}


%25.
\asana{[Lying] pronely, [the yogi] should put the toes on the ground, keep [the legs] long, place the palms of the hands at the top of the head, raise up the buttocks [and] gaze at the navel. [Then, the yogi] should put the nose onto the ground and take it as far as the palms of his hands. He should do this again and again. [This] is the \engname{Elephant Pose} (\engiast{gajāsana}).}
{ಹೊಟ್ಟೆ ಕೆಳಗೆ ಮಾಡಿ ಮಲಗಿ, ಕಾಲ್ಬೆರಳುಗಳನ್ನು ನೆಲಕ್ಕೆ ಊರಿ ಕಾಲುಗಳನ್ನು ಚಾಚಬೇಕು. ಅಂಗೈಗಳನ್ನು ತಲೆಯ ಮೇಲ್ಭಾಗದಲ್ಲಿ ಇಟ್ಟು, ಪೃಷ್ಠವನ್ನು ಮೇಲೆತ್ತಿ ಹೊಕ್ಕುಳನ್ನು ನೋಡಬೇಕು. ನಂತರ ಮೂಗನ್ನು ನೆಲಕ್ಕೆ ತಾಗಿಸಿ ಅಂಗೈಗಳವರೆಗೆ ತರಬೇಕು. ಇದನ್ನು ಪದೇ ಪದೇ ಮಾಡಬೇಕು. ಇದು \tbf{ಆನೆಯ ಭಂಗಿ} (\tit{ಗಜಾಸನ}).}
\end{paracol}
\vskip 0.8cm

\columnratio{0.45, 0.55}
\begin{paracol}{2}
\vfill
\begin{sans}
gajāsanavatsthitvā mastakaṃ vāraṃ vāraṃ dakṣiṇakukṣiṃ savyakukṣiṃ nayet\!। tarakṣvāsanaṃ bhavati\!॥26\!॥
\end{sans}

%26.
\asana{Remaining exactly as in the Elephant Pose, [the yogi] should take his head to the right armpit, [and then] to the left armpit, again and again. [This] is the \engname{Hyena Pose} (\engiast{tarakṣvāsana}).}
{ಗಜಾಸನದಲ್ಲಿರುವಂತೆಯೇ ಇದ್ದು, ತಲೆಯನ್ನು ಸರದಿಯಂತೆ ಬಲ ಮತ್ತು ಎಡ ಕಂಕುಳಿನ ಕಡೆಗೆ ಪದೇ ಪದೇ ತಿರುಗಿಸಬೇಕು. ಇದನ್ನು \tbf{ಕತ್ತೆಕಿರುಬನ ಭಂಗಿ} ಎನ್ನಲಾಗುತ್ತದೆ (\tit{ತರಕ್ಷ್ವಾಸನ}).}
\vfill
\switchcolumn

\insimgcol{HRCP/029_folio_08_verso_26.jpg}

\asananamestmp{tarakṣvāsanaṃ}{Hyena Pose}{ತರಕ್ಷ್ವಾಸನ}
\end{paracol}

\newpage

\columnratio{0.5, 0.5}
\begin{paracol}{2}
\insimgcol{HRCP/030_folio_09_recto_27.jpg}

\asananamestmp{ṛkṣāsanaṃ}{Bear Pose}{ಋಕ್ಷಾಸನ}
\bigskip

\begin{sans}
ekaikaṃ pādamākuñcya gajāsanaṃ kuryāt\!। ṛkṣāsanaṃ bhavati\!॥27\!॥
\end{sans}

%27.
\asana{[The yogi] should do the elephant pose, bending one leg at a time. [This] is the \engname{Bear Pose} (\engiast{ṛkṣāsana}).}
{ಗಜಾಸನದಲ್ಲಿದ್ದಾಗ ಒಂದೊಂದೇ ಕಾಲನ್ನು ಮಡಚುತ್ತಾ ಮಾಡುವ ಅಭ್ಯಾಸಕ್ಕೆ  \tbf{ಕರಡಿಯ ಭಂಗಿ} ಎನ್ನಲಾಗುತ್ತದೆ (\tit{ಋಕ್ಷಾಸನ}).}
\switchcolumn

%~ \end{paracol}
%~ \vskip 0.8cm

%~ \columnratio{0.5, 0.5}
%~ \begin{paracol}{2}
~\vskip 5cm
\begin{sans}
gajāsanasaṃsthitau\endnote{°saṃsthitau~] P, °sthitau M.} jānudvayamākuñcya vāraṃ vāraṃ kartavyam\!। śaśāsanaṃ bhavati\!॥28\!॥
\end{sans}

%28.
\asana{While in the Elephant Pose, [the yogi] should bend the knees again and again. [This] is the \engname{Hare Pose} (\engiast{śaśāsana}).}
{ಗಜಾಸನದಲ್ಲಿದ್ದಾಗ ಮೊಣಕಾಲುಗಳನ್ನು ಪದೇ ಪದೇ ಮಡಚುವುದಕ್ಕೆ \tbf{ಮೊಲದ ಭಂಗಿ} ಎನ್ನಲಾಗುತ್ತದೆ (\tit{ಶಶಾಸನ}).}
\bigskip

\insimgcol{HRCP/031_folio_09_recto_28.jpg}

\asananamestmp{śaśāsanaṃ}{Hare Pose}{ಶಶಾಸನ}

\end{paracol}

\newpage

\columnratio{0.5, 0.5}
\begin{paracol}{2}
~\vskip 4.5cm
\begin{sans}
gajāsanasaṃsthitau ekaikaṃ pādaṃ purobhāgena bhrāmayitvā kartavyam\!। rathāsanaṃ bhavati\!॥29\!॥
\end{sans}

%29.
\asana{While in the Elephant Pose, [the yogi] should rotate one leg at a time anteriorly and continue to do thus. [This] is the \engname{Chariot Pose} (\engiast{rathāsana}).}
{ಗಜಾಸನದಲ್ಲಿದ್ದಾಗ ಒಂದೊಂದೇ ಕಾಲನ್ನು ಮುಂದಕ್ಕೆ ವೃತ್ತಾಕಾರದಲ್ಲಿ ತಿರುಗಿಸುವುದಕ್ಕೆ \tbf{ತೇರಿನ ಭಂಗಿ} ಎನ್ನಲಾಗುತ್ತದೆ (\tit{ರಥಾಸನ}).}
\bigskip

\insimgcol{HRCP/032_folio_09_verso_29.jpg}

\asananamestmp{rathāsanaṃ}{Chariot Pose}{ರಥಾಸನ}



\switchcolumn
%~ \end{paracol}
%~ \vskip 0.8cm

%~ \columnratio{0.5, 0.5}
%~ \begin{paracol}{2}

\insimgcol{HRCP/033_folio_09_verso_30.jpg}

\asananamestmp{meṣāsanaṃ}{Sheep Pose}{ಮೇಷಾಸನ}
\bigskip

\begin{sans}
gajāsanasaṃsthitau ekaikaṃ bāhuṃ bhūmau tāḍayet\!। meṣāsanaṃ bhavati\!॥30\!॥
\end{sans}

%30.
\asana{While in the Elephant Pose, [the yogi] should hit the ground with one arm at a time. [This] is the \engname{Sheep Pose} (\engiast{meṣāsana}).}
{ಗಜಾಸನದಲ್ಲಿದ್ದಾಗ ಒಂದೊಂದೇ ತೋಳಿನಿಂದ ನೆಲಕ್ಕೆ ಹೊಡೆಯುವುದಕ್ಕೆ \tbf{ಕುರಿಯ ಭಂಗಿ} ಎನ್ನಲಾಗುತ್ತದೆ (\tit{ಮೇಷಾಸನ}).}


\end{paracol}

\newpage

\insimg[0.8]{HRCP/034_folio_10_recto_31.jpg}

\asananamestmpscale{ajāsanaṃ}{Goat Pose}{ಅಜಾಸನ}{0.8}
\vspace{0.5cm}

\begin{adjustwidth}{2.5cm}{2.5cm}
\begin{sans}
gajāsanasaṃsthitau pādadvayamantarālekṛtya mastakena bhūmiṃ spṛśet\!। ajāsanaṃ bhavati\!॥31\!॥
\end{sans}

%31.
\begin{multicols}{2}
\asana{While in the Elephant Pose, [the yogi] should raise both legs into the air and  touch the ground with the head. [This] is the \engname{Goat Pose} (\engiast{ajāsana}).}
{ಗಜಾಸನದಲ್ಲಿದ್ದಾಗ ಎರಡೂ ಕಾಲುಗಳನ್ನು ಮೇಲೆತ್ತಿ, ತಲೆಯಿಂದ ನೆಲವನ್ನು ಮುಟ್ಟಬೇಕು. ಇದನ್ನು \tbf{ಮೇಕೆಯ ಭಂಗಿ} ಎನ್ನಲಾಗುತ್ತದೆ (\tit{ಅಜಾಸನ}).}
\end{multicols}
\end{adjustwidth}

\newpage

\insimg[0.88]{HRCP/035_folio_10_recto_32.jpg}

\asananamestmpscale{caṭakāsanaṃ}{Sparrow Pose}{ಚಟಕಾಸನ}{0.88}
\vspace{0.5cm}

\begin{adjustwidth}{1.5cm}{1.5cm}
\begin{sans}
kūrparaparyantau hastau dharāmavaṣṭabhya jānunī nābhau saṃkuñcya  tiṣṭhet\!। caṭakāsanaṃ bhavati\!॥32\!॥
\end{sans}

%32.
\begin{multicols}{2}
\asana{Having supported [himself] with the forearms on the ground and bending the knees into the navel, [the yogi] should remain thus. [This] is the \engname{Sparrow Pose} (\engiast{caṭakāsana}).}
{ಮುಂಗೈಗಳನ್ನು ನೆಲಕ್ಕೆ ಊರಿ, ಮೊಣಕಾಲುಗಳನ್ನು ಮಡಚಿ ಹೊಕ್ಕುಳಿನ ಹತ್ತಿರ ತಂದು ಹಾಗೆಯೇ ಇರಬೇಕು. ಇದು \tbf{ಗುಬ್ಬಚ್ಚಿಯ ಭಂಗಿ} (\tit{ಚಟಕಾಸನ}).}
\end{multicols}
\end{adjustwidth}

\newpage

\columnratio{0.55, 0.45}
\begin{paracol}{2}

\insimgcol{HRCP/036_folio_10_verso_33.jpg}

\asananamestmp{kākāsanaṃ}{Crow Pose}{ಕಾಕಾಸನ}
\switchcolumn

\begin{sans}
hastau caṭakāsanavatsaṃsthāpya jānudvayaṃ karṇau saṃspṛśya jaṅghādvayaṃ bāhvornidhāya tiṣṭhet\!। kākāsanaṃ bhavati\!॥33\!॥
\end{sans}

%33.
\asana{Positioning the hands like the Sparrow Pose, [the yogi] should touch the ears with the knees, place the lower legs on the [upper] arms and remain thus. [This] is the \engname{Crow Pose} (\engiast{kākāsana}).}
{ಚಟಕಾಸನದಂತೆಯೇ ಕೈಗಳನ್ನು ಇಟ್ಟು, ಮೊಣಕಾಲುಗಳಿಂದ ಕಿವಿಗಳನ್ನು ಮುಟ್ಟಬೇಕು. ಕಾಲುಗಳ ಕೆಳಭಾಗವನ್ನು ತೋಳುಗಳ ಮೇಲಿಟ್ಟು ಹಾಗೆಯೇ ಇರಬೇಕು. ಇದನ್ನು  \tbf{ಕಾಗೆಯ ಭಂಗಿ} ಎನ್ನಲಾಗುತ್ತದೆ (\tit{ಕಾಕಾಸನ}).}
\end{paracol}
\vskip 0.4cm

\columnratio{0.45, 0.55}
\begin{paracol}{2}
\vfill
\begin{sans}
kākāsanavatsaṃsthitau jaṅghā ūruṇī saṃmīlya pṛṣṭhapradeśe ūrdhvaṃ nayet\!। tittirāsanaṃ bhavati\!॥34\!॥
\end{sans}

%34.
\asana{Positioned as in the Crow Pose, [the yogi] should join the shanks and thighs, and raise them up towards the back region [of his body]. [This] is the \engname{Partridge Pose} (\engiast{tittirāsana}).}
{ಕಾಕಾಸನದಲ್ಲಿದ್ದಾಗ ತೊಡೆ ಮತ್ತು ಕಾಲುಗಳನ್ನು ಜೋಡಿಸಿ ಬೆನ್ನಿನ ಕಡೆಗೆ ಮೇಲೆತ್ತಬೇಕು. ಇದು \tbf{ಗೌಜಲಕ್ಕಿಯ ಭಂಗಿ} (\tit{ತಿತ್ತಿರ್ಯಾಸನ}).}
\vfill
\switchcolumn
\insimgcol{HRCP/037_folio_10_verso_34.jpg}

\asananamestmp{tittirāsanaṃ}{Partridge Pose}{ತಿತ್ತಿರ್ಯಾಸನ}
\end{paracol}

\newpage

\columnratio{0.55, 0.45}
\begin{paracol}{2}
\insimgcol{HRCP/038_folio_11_recto_35.jpg}

\asananamestmp{bakāsanaṃ}{Heron Pose}{ಬಕಾಸನ}
\switchcolumn

\begin{sans}
hastābhyāmavanimavaṣṭabhya jānudvayaṃ nābhau saṃmīlya jaṅghā{}ūruṇī saṃstabhya tiṣṭhet\!। bakāsanaṃ bhavati\!॥35\!॥
\end{sans}

%35.
\asana{Supporting [himself] with the hands on the ground, [the yogi] should join the knees on the navel and hold the lower legs and thighs [in the air]. He should remain thus. [This] is the \engname{Heron Pose} (\engiast{bakāsana}).}
{ಅಂಗೈಗಳನ್ನು ನೆಲಕ್ಕೆ ಊರಿ, ಮೊಣಕಾಲುಗಳನ್ನು ಹೊಕ್ಕುಳಿನ ಹತ್ತಿರ ಜೋಡಿಸಿ, ಕಾಲುಗಳನ್ನು ಗಾಳಿಯಲ್ಲಿ ಹಿಡಿಯಬೇಕು. ಇದನ್ನು \tbf{ಕೊಕ್ಕರೆಯ ಭಂಗಿ} ಎನ್ನಲಾಗುತ್ತದೆ (\tit{ಬಕಾಸನ}).}
\end{paracol}
\vskip 0.5cm

\columnratio{0.45, 0.55}
\begin{paracol}{2}
\vfill
\begin{sans}
padmāsanaṃ saṃsthāpya hastatalābhyāṃ dharāmavaṣṭabhya padmāsanayuktacaraṇadvayaṃ pṛṣṭhabhāge nītvā tiṣṭhet\!। bhāradvājāsanaṃ bhavati\!॥36\!॥
\end{sans}

%36.
\asana{Having adopted the Lotus Pose, [the yogi] should support [himself] with the palms of the hands on the ground, lift the feet, [which are] fastened in lotus pose, towards the back and remain thus. [This] is \engname{Bhāradvāja’s Pose} (\engiast{bhāradvājāsana}).
}{ಪದ್ಮಾಸನ ಹಾಕಿಕೊಂಡು ಅಂಗೈಗಳನ್ನು ನೆಲಕ್ಕೆ ಊರಿ, ಪದ್ಮಾಸನದಲ್ಲಿರುವ ಕಾಲುಗಳನ್ನು ಬೆನ್ನಿನ ಕಡೆಗೆ ಮೇಲೆತ್ತಬೇಕು. ಇದನ್ನು \tbf{ಭಾರದ್ವಾಜಾಸನ} ಎನ್ನಲಾಗುತ್ತದೆ.}
\vfill
\switchcolumn

\insimgcol{HRCP/039_folio_11_recto_36.jpg}

\asananamestmp{bhāradvājāsanaṃ}{Bhāradvāja’s Pose}{ಭಾರದ್ವಾಜಾಸನ}
\end{paracol}

\newpage

\insimg[0.9]{HRCP/040_folio_11_verso_37.jpg}

\asananamestmpscale{kukkuṭṭoḍḍānaṃ}{Flight of the Cock [Pose]}{ಕುಕ್ಕುಟೋಡ್ಡಯನಾಸನ}{0.9}
\vspace{0.5cm}

\begin{adjustwidth}{1.3cm}{1.3cm}
\begin{sans}
hastatale bhūmau kṛtvā pādatale ca ūrdhvamuḍḍānaṃ kṛtvā bhūmau patet\!। itthaṃ punaḥ punaḥ kuryāt\!। kukkuṭṭoḍḍānaṃ bhavati\!॥37\!॥
\end{sans}

%37.
\begin{multicols}{2}
\asana{Putting the palms of the hands on the ground, [the yogi] should make the soles of the feet fly upwards and then fall down to the ground. He should do thus again and again. [This] is the \engname{Flight of the Cock [Pose]} (\engiast{kukkuṭṭoḍḍāna}).}
{ಅಂಗೈಗಳನ್ನು ನೆಲಕ್ಕೆ ಊರಿ, ಪಾದಗಳನ್ನು ಗಾಳಿಯಲ್ಲಿ ಮೇಲಕ್ಕೆ ಹಾರಿಸಿ ಮತ್ತೆ ನೆಲಕ್ಕೆ ತಾಗಿಸಬೇಕು. ಇದನ್ನು ಪದೇ ಪದೇ ಮಾಡಬೇಕು. ಇದು \tbf{ಹುಂಜದ ಹಾರಾಟದ ಭಂಗಿ} (\tit{ಕುಕ್ಕುಟೋಡ್ಡಯನಾಸನ}).}
\end{multicols}
\end{adjustwidth}

\newpage

\insimg[0.9]{HRCP/041_folio_11_verso_38.jpg}

\asananamestmpscale{araṇyacaṭakāsanaṃ}{Wood-sparrow Pose}{ಅರಣ್ಯಚಟಕಾಸನ}{0.9}
\vspace{0.5cm}

\begin{adjustwidth}{1.3cm}{1.3cm}
\begin{sans}
ekaṃ pādaṃ grīvāyāṃ saṃsthāpya dvitīyaṃ pādaṃ upari vinyasya hastatalābhyāmavaṣṭabhya tiṣṭhet\!। araṇyacaṭakāsanaṃ bhavati\!॥38\!॥
\end{sans}

%38.
\begin{multicols}{2}
\asana{Having placed one foot on [the back of] the neck, [the yogi] should fix the other foot above it, support [the body] with the palms of the hands [on the ground] and remain thus. [This] is the \engname{Wood-sparrow Pose} (\engiast{araṇyacaṭakāsana}).}
{ಒಂದು ಪಾದವನ್ನು ಕುತ್ತಿಗೆಯ ಹಿಂದೆ ಇಟ್ಟು, ಅದರ ಮೇಲೆ ಇನ್ನೊಂದು ಪಾದವನ್ನು ಇರಿಸಬೇಕು. ಅಂಗೈಗಳಿಂದ ಶರೀರವನ್ನು ನೆಲದ ಮೇಲೆ ಸಮತೋಲನಗೊಳಿಸಬೇಕು. ಇದು \tbf{ಮರಗುಬ್ಬಚ್ಚಿಯ ಭಂಗಿ} (\tit{ಅರಣ್ಯಚಟಕಾಸನ}).}
\end{multicols}
\end{adjustwidth}
\newpage

\columnratio{0.5, 0.5}
\begin{paracol}{2}
\insimgcol{HRCP/042_folio_12_recto_39.jpg}

\asananamestmp{mayūrāsanaṃ}{Peacock Pose}{ಮಯೂರಾಸನ}
\bigskip

\begin{sans}
hastatalābhyāmavanimavaṣṭabhya kūrparau nābhau vinyasya daṇḍavaccharīraṃ dhṛtvā tiṣṭhet\!। mayūrāsanaṃ bhavati\!॥39\!॥
\end{sans}

%39.
\asana{Supporting [himself] with the palms of the hands on the ground, fixing the elbows on the navel and holding the body [straight] like a stick, [the yogi] should remain thus. [This] is the \engname{Peacock Pose} (\engiast{mayūrāsana}).}
{ಅಂಗೈಗಳನ್ನು ನೆಲಕ್ಕೆ ಊರಿ, ಮೊಣಕೈಗಳನ್ನು ಹೊಕ್ಕುಳಿನ ಮೇಲೆ ಸ್ಥಿರವಾಗಿರಿಸಿ, ಇಡೀ ಶರೀರವನ್ನು ಕೋಲಿನಂತೆ ನೇರವಾಗಿ ಗಾಳಿಯಲ್ಲಿ ಹಿಡಿಯಬೇಕು. ಇದು \tbf{ನವಿಲಿನ ಭಂಗಿ} (\tit{ಮಯೂರಾಸನ}).}
\switchcolumn
~\vskip 4cm
%~ \end{paracol}
%~ \vskip 0.4cm

%~ \columnratio{0.5, 0.5}
%~ \begin{paracol}{2}
\begin{sans}
mayūrāsanavatsaṃsthāpya ekena hastena ekasya hastasya maṇibandhaṃ dhārayet\!। paṅgumayūraṃ bhavati\!॥40\!॥
\end{sans}

%40.
\asana{Having positioned [himself] as in the Peacock Pose, [the yogi] should hold the wrist of one hand with the other hand. [This] is the \engname{Lame Peacock [Pose]} (\engiast{paṅgumayūra}).}
{ಮಯೂರಾಸನದಲ್ಲಿರುವಾಗ ಒಂದು ಕೈಯಿಂದ ಇನ್ನೊಂದು ಕೈಯ ಮಣಿಕಟ್ಟನ್ನು ಹಿಡಿಯಬೇಕು. ಇದನ್ನು \tbf{ಕುಂಟು ನವಿಲಿನ ಭಂಗಿ} ಎನ್ನಲಾಗುತ್ತದೆ (\tit{ಪಂಗುಮಯೂರಾಸನ}).}
\bigskip

\insimgcol{HRCP/043_folio_12_recto_40.jpg}

\asananamestmp{paṅgumayūraṃ}{Lame Peacock Pose}{ಪಂಗುಮಯೂರಾಸನ}
\end{paracol}


\newpage

\insimg[0.76]{HRCP/044_folio_12_verso_41.jpg}

\asananamestmpscale{khaḍgāsanaṃ}{Sword Pose}{ಖಡ್ಗಾಸನ}{0.76}
\vspace{0.5cm}

\begin{adjustwidth}{2.8cm}{2.8cm}
\begin{sans}
nyubjaṃ saralamāsanaṃ kṛtvā pādatalābhyāṃ bhūmimavaṣṭabhya uttiṣṭhet\!। khaḍgāsanaṃ bhavati\!॥41\!॥
\end{sans}

%41
\begin{multicols}{2}
\asana{[The yogi] should lie prone, sit up straight, support [himself ] with the soles of the feet on the
ground, and stand up. [This] is the \engname{Sword Pose} (\engiast{khaḍgāsana}).}
{ಹೊಟ್ಟೆ ಕೆಳಗೆ ಮಾಡಿ ಮಲಗಿ, ನಂತರ ನೇರವಾಗಿ ಕುಳಿತು, ಪಾದಗಳನ್ನು ನೆಲಕ್ಕೆ ಊರಿ ಎದ್ದು ನಿಲ್ಲಬೇಕು. ಇದು \tbf{ಕತ್ತಿಯ ಭಂಗಿ} (\tit{ಖಡ್ಗಾಸನ}).}
\end{multicols}
\end{adjustwidth}
\newpage

\insimg[0.9]{HRCP/045_folio_12_verso_42.jpg}

\asananamestmpscale{śūlāsanaṃ}{Spear Pose}{ಶೂಲಾಸನ}{0.9}
\vspace{0.5cm}

\begin{adjustwidth}{1.3cm}{1.3cm}
\begin{sans}
kūrparābhyāmavanimavaṣṭabhya hastatalābhyāṃ hanumavaṣṭabhya uttiṣṭhet\!। śūlāsanaṃ bhavati\!॥42\!॥
\end{sans}

%42
\begin{multicols}{2}
\asana{Supporting [himself ] with the elbows on the ground and supporting the jaw with the palms of the
hands, [the yogi] should raise himself up. [This] is the \engname{Spear Pose} (\engiast{śūlāsana}).}
{ಮೊಣಕೈಗಳನ್ನು ನೆಲಕ್ಕೆ ಊರಿ, ಗಲ್ಲವನ್ನು ಅಂಗೈಗಳ ಮೇಲೆ ಇಟ್ಟುಕೊಂಡು ಶರೀರವನ್ನು ಮೇಲೆತ್ತಬೇಕು. ಇದು \tbf{ಈಟಿಯ ಭಂಗಿ}. (\tit{ಶೂಲಾಸನ}).}
\end{multicols}
\end{adjustwidth}
\newpage

\insimg[0.81]{HRCP/046_folio_13_recto_43.jpg}

\asananamestmpscale{viparītanṛtyaṃ}{Inverted Dancing [Pose]}{ವಿಪರೀತನೃತ್ಯಾಸನ}{0.81}
\vspace{0.5cm}

\begin{adjustwidth}{2.4cm}{2.4cm}
\begin{sans}
hastatalābhyāṃ bhūmimavaṣṭabhya ūrdhvaṃ pādāgre kurvanhastatalābhyāṃ nartanaṃ kartavyam\!। viparītanṛtyaṃ bhavati\!॥43\!॥ 
\end{sans}

%43
\begin{multicols}{2}
\asana{Having supported [himself ] with the palms of the hands on the ground and lifting the toes up
[into the air], [the yogi] should dance on the palms of the hands. [This] is the \engname{Inverted Dancing
[Pose]} (\engiast{viparītanṛtya}).}
{ಅಂಗೈಗಳನ್ನು ನೆಲಕ್ಕೆ ಊರಿ, ಕಾಲ್ಬೆರಳುಗಳನ್ನು ಗಾಳಿಯಲ್ಲಿ ಮೇಲೆತ್ತಿ ಅಂಗೈಗಳ ಮೇಲೆ ನೃತ್ಯ ಮಾಡುವಂತೆ ಚಲಿಸಬೇಕು. ಇದನ್ನು \tbf{ವಿಪರೀತನೃತ್ಯಾಸನ} ಎನ್ನಲಾಗುತ್ತದೆ.}
\end{multicols}
\end{adjustwidth}

\newpage

\insimg[0.83]{HRCP/047_folio_13_recto_44.jpg}

\asananamestmpscale{śyenāsanaṃ}{Hawk Pose}{ಶ್ಯೇನಾಸನ}{0.83}
\vspace{0.5cm}

\begin{adjustwidth}{2cm}{2cm}
\begin{sans}
viparītanṛtyavatsaṃsthitau nāsikāṃ\endnote{nāsikāṃ~] \textit{corr.},  nāśikāṃ M.} bhūmau saṃspṛśya ūrdhvaṃ nītvā punaḥ saṃspṛśet\!। śyenāsanaṃ bhavati\!॥44\!॥
\end{sans}

%44
\begin{multicols}{2}
\asana{Positioned as in the Inverted Dancing Pose, [the yogi] should touch the nose on the ground, lift it
up and touch [the ground] again [and again]. [This] is the \engname{Hawk Pose} (\engiast{śyenāsana}).}
{ವಿಪರೀತ ನೃತ್ಯಾಸನದಲ್ಲಿದ್ದಾಗ ಮೂಗನ್ನು ನೆಲಕ್ಕೆ ತಾಗಿಸಿ ಮತ್ತೆ ಮೇಲೆತ್ತಬೇಕು. ಇದನ್ನು \tbf{ಗಿಡುಗದ ಭಂಗಿ} ಎನ್ನಲಾಗುತ್ತದೆ (\tit{ಶ್ಯೇನಾಸನ}).}
\end{multicols}
\end{adjustwidth}

\newpage


\insimg[0.7]{HRCP/048_folio_13_verso_45.jpg}

\asananamestmpscale{kapālāsanaṃ}{Skull Pose}{ಕಪಾಲಾಸನ}{0.7}
\vspace{0.5cm}

\begin{adjustwidth}{3.65cm}{3.65cm}
\begin{sans}
kapālaṃ bhūmau nidhāya ūrdhvaṃ pādau nayet\!। kapālāsanaṃ bhavati\!॥45\!॥
\end{sans}

%45
\begin{multicols}{2}
\asana{Placing the [top of the] skull on the ground, [the yogi] should lift the feet [into the air]. [This] is
the \engname{Skull Pose} (\engiast{kapālāsana}).}
{ತಲೆಯ ಮೇಲ್ಭಾಗವನ್ನು (ಕಪಾಲ) ನೆಲಕ್ಕೆ ಊರಿ, ಕಾಲುಗಳನ್ನು ಗಾಳಿಯಲ್ಲಿ ಮೇಲೆತ್ತಬೇಕು. ಇದು \tbf{ಕಪಾಲಾಸನ} (ಶೀರ್ಷಾಸನದ ಒಂದು ರೂಪ).}
\end{multicols}
\end{adjustwidth}


\newpage

\columnratio{0.5, 0.5}
\begin{paracol}{2}
\insimgcol{HRCP/049_folio_13_verso_46.jpg}

\asananamestmp{sarpāsanaṃ}{Snake Pose}{ಸರ್ಪಾಸನ}
\bigskip

\begin{sans}
nyubjaṃ śayanaṃ kṛtvā hastau nitambe saṃsthāpya pādau dīrghīkṛtya saṃmīlayannurasā\endnote{mīlayann~] \textit{corr.}, mīlayan M.} gantavyam\!। sarpāsanaṃ bhavati\!॥46\!॥
\end{sans}

%46
\asana{Lying prone, [the yogi] should place the hands on the buttocks, extend the legs, join [them]
together, and move [forwards] with his chest. [This] is the \engname{Snake Pose} (\engiast{sarpāsana}).}{ಹೊಟ್ಟೆ ಕೆಳಗೆ ಮಾಡಿ ಮಲಗಿ, ಕೈಗಳನ್ನು ಪೃಷ್ಠದ ಮೇಲಿಟ್ಟು, ಕಾಲುಗಳನ್ನು ಚಾಚಿ ಜೋಡಿಸಿ ಎದೆಯ ಭಾಗದಿಂದ ಮುಂದಕ್ಕೆ ಸರಿಯಬೇಕು. ಇದು \tbf{ಹಾವಿನ ಭಂಗಿ} (\tit{ಸರ್ಪಾಸನ}).}
\switchcolumn
~\vskip 3cm

\begin{sans}
nyubjaṃ śayanaṃ kṛtvā pṛṣṭhe pārṣṇivyutkrameṇa\endnote{pārṣṇi°~] \textit{corr.}~: pārṣṇī M.} hastābhyāṃ  pādāṅguṣṭhau dhṛtvā loḍayet\!। nyubjāsanaṃ bhavati\endnote{nyubjāsanaṃ bhavati~] \textit{diagnostic conj.}~: \textit{omitted.} M.}\!॥47\!॥
\end{sans}

%47
\asana{[The yogi] should lie prone, cross the heels behind him, hold the big toes with the hands, and roll
[around]. [This] is the \engname{Prone Pose} (\engiast{nyubjāsana}).}{ಹೊಟ್ಟೆ ಕೆಳಗೆ ಮಾಡಿ ಮಲಗಿ, ಹಿಮ್ಮಡಿಗಳನ್ನು ಹಿಂದಕ್ಕೆ ಕತ್ತರಿ ಆಕಾರದಲ್ಲಿ ಜೋಡಿಸಿ, ಕೈಗಳಿಂದ ಕಾಲ್ಬೆರಳುಗಳನ್ನು ಹಿಡಿದು ಉರುಳಬೇಕು. ಇದನ್ನು  \tbf{ನ್ಯುಬ್ಜಾಸನ} ಎನ್ನಲಾಗುತ್ತದೆ.}
\bigskip

\insimgcol{HRCP/050_folio_14_recto_47_E_PRONE.jpg}

\asananamestmp{nyubjāsanaṃ}{Prone Pose}{ನ್ಯುಬ್ಜಾಸನ}
\end{paracol}
\newpage

\insertFullImage{057_folio_15_verso_54.jpg}

\newpage

\pagesWithBorder

\null\vfill
{\leftskip=30pt
\begin{sans}
{\fontsize{18pt}{21.6pt}\selectfont \color{deepmaroon}{atha sthānāsanyāsanāni\!॥}}
\end{sans}
\vskip 7pt

\begin{eng}
\prtheading[Stationary Poses]{Now, the Stationary Poses}
\end{eng}
\vskip 7pt

\begin{kan}
\prtheading[ಕುಳಿತು ಮಾಡುವ ಅಥವಾ ಸ್ಥಿರವಾಗಿರುವ ಆಸನಗಳು]{ಈಗ, ಕುಳಿತು ಮಾಡುವ ಅಥವಾ ಸ್ಥಿರವಾಗಿರುವ ಆಸನಗಳು}
\end{kan}\par}
\vfill

\newpage

\insimg[0.82]{HRCP/051_folio_14_recto_48_B_STATIONARY.jpg}

\asananamestmpscale{paścimatānāsanaṃ}{Back Stretch Pose}{ಪಶ್ಚಿಮತಾನಾಸನ}{0.82}
\vspace{0.5cm}

\begin{adjustwidth}{2.3cm}{2.3cm}
\begin{sans}
daṇḍavadbhūmau caraṇau prasārya hastābhyāmaṅguṣṭhau dhṛtvā jānūpari lalāṭaṃ vinyasya tiṣṭhet\!। paścimatānāsanaṃ bhavati\endnote{paścimatānāsanaṃ bhavati~] \textit{diagnostic conj.}~: \textit{omitted} M.}\!॥48\!॥
\end{sans}

%48
\begin{multicols}{2}
\asana{[The yogi] should extend the legs on the ground like a stick, hold the big toes with the hands,
place the forehead on the knees and remain thus. [This] is the \engname{Back Stretch Pose}
(paścimatānāsana).}{ಕಾಲುಗಳನ್ನು ನೆಲದ ಮೇಲೆ ಕೋಲಿನಂತೆ ಚಾಚಿ, ಕೈಗಳಿಂದ ಕಾಲಿನ ಹೆಬ್ಬೆರಳುಗಳನ್ನು ಹಿಡಿದು, ಹಣೆಯನ್ನು ಮೊಣಕಾಲುಗಳಿಗೆ ತಾಗಿಸಿ ಹಾಗೆಯೇ ಇರಬೇಕು. ಇದು  \tbf{ಬೆನ್ನು ಸೆಳೆಯುವ ಭಂಗಿ} (\tit{ಪಶ್ಚಿಮತಾನಾಸನ}).}
\end{multicols}
\end{adjustwidth}

\newpage

\columnratio{0.5, 0.5}
\begin{paracol}{2}
\insimgcol{HRCP/052_folio_14_verso_49.jpg}

\asananamestmp{ardhapaścimatānaṃ}{Half Back Stretch [Pose]}{ಅರ್ಧಪಶ್ಚಿಮತಾನಾಸನ}
\switchcolumn

\begin{sans}
ekaṃ pādaṃ prasārya ekapādasya pārṣṇinā sīvanīṃ saṃpīḍya hastadvayena prasāritapādāṅguṣṭhaṃ dhṛtvā jānūpari mastakaṃ nyaset\!। ardhapaścimatānaṃ bhavati\!॥49\!॥
\end{sans}

%49
\asana{[The yogi] should extend one leg, press the perineal seam with the heel of the other foot, hold the
big toe of the extended leg with the hands and fix the head on the knee. [This] is the \engname{Half Back Stretch [Pose]} (\engiast{ardhapaścimatāna}).}{ಒಂದು ಕಾಲನ್ನು ಚಾಚಿ, ಇನ್ನೊಂದು ಕಾಲಿನ ಹಿಮ್ಮಡಿಯಿಂದ ಶಿಶ್ನ ಮತ್ತು ಗುದದ್ವಾರದ ಮಧ್ಯದ ಭಾಗವನ್ನು ಒತ್ತಿ ಹಿಡಿಯಬೇಕು. ಚಾಚಿದ ಕಾಲಿನ ಹೆಬ್ಬೆರಳನ್ನು ಕೈಗಳಿಂದ ಹಿಡಿದು ತಲೆಯನ್ನು ಮೊಣಕಾಲಿಗೆ ತಾಗಿಸಬೇಕು. ಇದು \tbf{ಅರ್ಧಪಶ್ಚಿಮತಾನಾಸನ}.}
\end{paracol}
\vskip 0.8cm

\columnratio{0.5, 0.5}
\begin{paracol}{2}
\vfill
\begin{sans}
paścimatānavatsaṃsthitiṃ kṛtvā ekaṃ pādaṃ grīvāyāṃ dhṛtvā tiṣṭhet\!। ūrdhvapaścimatānaṃ bhavati\!॥50\!॥
\end{sans}

%50
\asana{Positioning [himself ] as in Stretching the Back [Pose] and holding one foot on [the back of ] the
neck, [the yogi] should remain thus. [This] is the \engname{Upper Back Stretch [Pose]} (\engiast{ūrdhvapaścimatāna}).}
{ಪಶ್ಚಿಮತಾನಾಸನದ ಸ್ಥಿತಿಯಲ್ಲಿದ್ದಾಗ ಒಂದು ಕಾಲನ್ನು ಕುತ್ತಿಗೆಯ ಹಿಂಭಾಗದಲ್ಲಿ ಇರಿಸಿಕೊಳ್ಳಬೇಕು. ಇದು \tbf{ಊರ್ಧ್ವಪಶ್ಚಿಮತಾನಾಸನ}.}
\vfill
\switchcolumn
\insimgcol{HRCP/053_folio_14_verso_50.jpg}

\asananamestmp{ūrdhvapaścimatānaṃ}{Upper Back Stretch [Pose]}{ಊರ್ಧ್ವಪಶ್ಚಿಮತಾನಾಸನ}
\end{paracol}

\newpage

\columnratio{0.5, 0.5}
\begin{paracol}{2}
\insimgcol{HRCP/054_folio_15_recto_51.jpg}

\asananamestmp{dhanurāsanaṃ}{Bow Pose}{ಧನುರಾಸನ}
\vskip 0.5cm
%~ \switchcolumn

\begin{sans}
hastadvayena pādadvayāgre gṛhītvā ekaikaṃ pādāṅguṣṭhaṃ karṇayoḥ spṛśet\!। dhanurāsanaṃ bhavati\endnote{bhavati~] \textit{conj.}, \textit{omitted.} M.}\!॥51\!॥
\end{sans}

%51
\asana{Having grasped the toes of the feet with the hands, [the yogi] should touch the big toes, one at a time, on the ears. This is the \engname{Bow Pose} (\engiast{dhanurāsana}).}
{ಕೈಗಳಿಂದ ಕಾಲ್ತುದಿಯನ್ನು ಹಿಡಿದು, ಒಂದೊಂದೇ ಕಾಲಿನ ಹೆಬ್ಬೆರಳನ್ನು ಕಿವಿಯವರೆಗೆ ತಂದು ಮುಟ್ಟಿಸಬೇಕು. ಇದು \tbf{ಬಿಲ್ಲಿನ ಭಂಗಿ} (\tit{ಧನುರಾಸನ}).}
\switchcolumn
~\vskip 4cm
%~ \end{paracol}
%~ \vskip 0.8cm

%~ \columnratio{0.5, 0.5}
%~ \begin{paracol}{2}
%~ \vfill
\begin{sans}
doḥkuṭṭanam\!। ūrukuṭṭanam\!। pārśvakuṭṭanam\!। ityādīni kuṭṭanāni muṣṭinā\!। bāhunā pārṣṇinā bhityā bhūmyā kartavyāni\!॥52\!॥
\end{sans}

%52
\asana{Hitting the arms, thighs, the sides and so on: the strikes should be done with a fist, arm, heel, wall [or] the ground.}
{ತೋಳುಗಳು, ತೊಡೆಗಳು ಮತ್ತು ಪಕ್ಕೆಲುಬುಗಳ ಭಾಗಕ್ಕೆ ಮುಷ್ಟಿಯಿಂದ ಅಥವಾ ತೋಳಿನಿಂದ ಪೆಟ್ಟು ನೀಡಬೇಕು. ಹಾಗೆಯೇ ಹಿಮ್ಮಡಿಯಿಂದ, ಗೋಡೆಗೆ ಅಥವಾ ನೆಲಕ್ಕೆ ದೇಹದ ಭಾಗಗಳನ್ನು ತಾಗಿಸುವ ಮೂಲಕ ಹೊಡೆತಗಳನ್ನು ನೀಡಬೇಕು. ಅಂದರೆ, ಶರೀರವನ್ನು ಬಲಪಡಿಸಲು ಮುಷ್ಟಿ, ಹಿಮ್ಮಡಿ, ಗೋಡೆ ಅಥವಾ ನೆಲದ ಸಹಾಯದಿಂದ ಮೈಮೇಲೆ ಪೆಟ್ಟುಗಳನ್ನು ಹಾಕಿಕೊಳ್ಳಬೇಕು.}
\vskip 0.5cm

\insimgcol{HRCP/055_folio_15_recto_52.jpg}
\end{paracol}

\newpage

\insimg[0.88]{HRCP/056_folio_15_verso_53.jpg}

\asananamestmpscale{baddhapadmāsanaṃ}{Bound Lotus Pose}{ಬದ್ಧ ಪದ್ಮಾಸನ}{0.88}
\vspace{0.5cm}

\begin{adjustwidth}{1.42cm}{1.42cm}
\begin{sans}
ūruṇi pādaṃ vinyasya itarasminnūruṇi itaraṃ pādaṃ nyasya pṛṣṭhabhāgena\endnote{pṛṣṭha°~] \textit{corr.}, pṛṣṭhi° M.} vyatyayena hastābhyāṃ pādāṅguṣṭhau dhṛtvā tiṣṭhet\!। baddhapadmāsanaṃ bhavati\!॥52\,[53]\!॥
\end{sans}

%53
\begin{multicols}{2}
\asana{[The yogi] should fix [one] foot on the [opposite] thigh, the other foot on the other thigh, and
hold the big toes with the hands crossed behind the back. He should remain thus. [This] is the
\engname{Bound Lotus Pose} (\engiast{baddhapadmāsana}).}{ಒಂದು ಪಾದವನ್ನು ಎದುರಿನ ತೊಡೆಯ ಮೇಲೂ, ಇನ್ನೊಂದು ಪಾದವನ್ನು ಮತ್ತೊಂದು ತೊಡೆಯ ಮೇಲೂ ಇರಿಸಿ (ಪದ್ಮಾಸನ), ಬೆನ್ನಿನ ಹಿಂದಿನಿಂದ ಕೈಗಳನ್ನು ಕತ್ತರಿ ಆಕಾರದಲ್ಲಿ ತಂದು ಕಾಲ್ಬೆರಳುಗಳನ್ನು ಹಿಡಿಯಬೇಕು. ಇದು \tbf{ಬದ್ಧ ಪದ್ಮಾಸನ}.}
\end{multicols}
\end{adjustwidth}

\newpage

\insimg[0.94]{HRCP/057_folio_15_verso_54.jpg}

\asananamestmpscale{kukkuṭāsanaṃ}{Cock Pose}{ಕುಕ್ಕುಟಾಸನ}{0.94}
\vspace{0.5cm}

\begin{adjustwidth}{0.8cm}{0.8cm}
\begin{sans}
padmāsanaṃ kṛtvā caraṇa{}ūrujaṅghānāmantare bāhudvayaṃ nyasya  hastatalābhyāṃ bhūmimavaṣṭabhya tiṣṭhet\!। kukkuṭāsanaṃ bhavati\!॥53\,[54]\!॥
\end{sans}

%54
\begin{multicols}{2}
\asana{Having adopted the Lotus Pose, [the yogi] should insert the arms between the feet, thighs and calves, support [himself ] with the palms of the hands on the ground, and remain thus. [This] is the \engname{Cock Pose} (\engiast{kukkuṭāsana}).}
{ಪದ್ಮಾಸನದಲ್ಲಿದ್ದಾಗ ಪಾದ ಮತ್ತು ತೊಡೆಗಳ ಮಧ್ಯದ ಸಂದಿನಿಂದ ಕೈಗಳನ್ನು ತೂರಿಸಿ, ಅಂಗೈಗಳನ್ನು ನೆಲಕ್ಕೆ ಊರಿ ಶರೀರವನ್ನು ಮೇಲೆತ್ತಬೇಕು. ಇದು \tbf{ಹುಂಜದ ಭಂಗಿ} (\tit{ಕುಕ್ಕುಟಾಸನ}).}
\end{multicols}
\end{adjustwidth}

\newpage


\columnratio{0.5, 0.5}
\begin{paracol}{2}
\insimgcol{HRCP/058_folio_16_recto_55.jpg}

\asananamestmp{paṅgukukkuṭaṃ}{Lame Cock [Pose]}{ಪಂಗುಕುಕ್ಕುಟಾಸನ}
\vskip 0.5cm

\begin{sans}
kukkuṭāsanavatsthitvā ekena hastena anyasya hastasya maṇibandhaṃ dhṛtvā hastatalena bhūmiṃ viṣṭabhya tiṣṭhet\!। paṅgukukkuṭaṃ bhavati\!॥54\,[55]\!॥
\end{sans}

%55
\asana{While in the Cock Pose, [the yogi] should hold the wrist of one hand with the other hand, support himself with the palm of the [held] hand on the ground and remain thus. [This] is the \engname{Lame Cock [Pose]} (\engiast{paṅgukukkuṭa}).}
{ಕುಕ್ಕುಟಾಸನದಲ್ಲಿದ್ದಾಗ (ಹುಂಜದ ಭಂಗಿ), ಯೋಗಿಯು ತನ್ನ ಒಂದು ಕೈಯ ಮಣಿಕಟ್ಟನ್ನು ಇನ್ನೊಂದು ಕೈಯಿಂದ ಗಟ್ಟಿಯಾಗಿ ಹಿಡಿಯಬೇಕು. ನಂತರ ಆ ಹಿಡಿದ ಕೈಯ ಅಂಗೈಯನ್ನು ಮಾತ್ರ ನೆಲಕ್ಕೆ ಊರಿ ಇಡೀ ಶರೀರವನ್ನು ಸಮತೋಲನಗೊಳಿಸಬೇಕು. ಇದನ್ನು \tbf{ಕುಂಟು ಹುಂಜದ ಭಂಗಿ} ಎಂದು ಕರೆಯಲಾಗುತ್ತದೆ (\tit{ಪಂಗುಕುಕ್ಕುಟಾಸನ}).}
\switchcolumn
%~ \end{paracol}
%~ \vskip 0.8cm

%~ \columnratio{0.5, 0.5}
%~ \begin{paracol}{2}
%~ \vfill
~\vskip 4cm
\begin{sans}
nyubjaṃ śayanaṃ kṛtvā pārṣṇidvayaṃ grīvāyāṃ sthāpayet hastadvayena gulphadvayaṃ dhṛtvā tiṣṭhet\!। arghyāsanaṃ bhavati\endnote{arghyāsanaṃ bhavati~] \textit{diagnostic conj.}~: \textit{omitted} M.}\!॥55\,[56]\!॥
\end{sans}

%56
\asana{Lying prone, [the yogi] should place the heels on the neck. Holding both ankles with the hands, he
should remain thus. This is the \engname{Libation-bowl Pose} (\engiast{arghyāsana}).}
{ಹೊಟ್ಟೆ ಕೆಳಗೆ ಮಾಡಿ ಮಲಗಿ ಹಿಮ್ಮಡಿಗಳನ್ನು ಕುತ್ತಿಗೆಯ ಮೇಲೆ ಇರಿಸಬೇಕು. ಕೈಗಳಿಂದ ಎರಡು ಪಾದದ ಗಂಟನ್ನು ಹಿಡಿದು ಹಾಗೆಯೇ ಇರಬೇಕು. ಇದು \tbf{ಅರ್ಘ್ಯಾಸನ}.}
%~ \vfill
%~ \switchcolumn
\bigskip

\insimgcol{HRCP/059_folio_16_recto_56.jpg}

\asananamestmp{arghyāsanaṃ}{Libation-bowl Pose}{ಅರ್ಘ್ಯಾಸನ}
\end{paracol}

\newpage

\columnratio{0.5, 0.5}
\begin{paracol}{2}
\insimgcol{HRCP/060_folio_16_verso_57.jpg}

\asananamestmp{chatrāsanaṃ}{Parasol Pose}{ಛತ್ರಾಸನ}
%~ \switchcolumn
\vskip 0.5cm

\begin{sans}
arghyāsanavatsthitvā hastatalābhyāṃ bhūmimavaṣṭabhnuyāt\!। chatrāsanaṃ bhavati\!॥56\,[57]\!॥
\end{sans}

%57
\asana{Positioned as in the Libation-bowl Pose, [the yogi] should support [himself ] with the palms of the
hands on the ground. [This] is the \engname{Parasol Pose} (\engiast{chatrāsana}).}{ಅರ್ಘ್ಯಾಸನದಲ್ಲಿದ್ದಾಗ ಅಂಗೈಗಳನ್ನು ನೆಲಕ್ಕೆ ಊರಿ ಸಮತೋಲನ ಕಾಪಾಡಿಕೊಳ್ಳಬೇಕು. ಇದು \tbf{ಛತ್ರಿಯ ಭಂಗಿ} (\tit{ಛತ್ರಾಸನ}).}
\switchcolumn
~\vskip 3.5cm
%~ \end{paracol}
%~ \vskip 0.8cm

%~ \columnratio{0.5, 0.5}
%~ \begin{paracol}{2}
%~ \vfill
\begin{sans}
hastābhyāmavanimavaṣṭabhya skandhayorjānunī\endnote{skandhayor~] \textit{corr.}~:  skandhayo M.} saṃsthāpya pārṣṇī\endnote{pārṣṇī~] \textit{corr.}~: pārṣṇi M.} urasi nidhāya tiṣṭhet\!। mālāsanaṃ bhavati\!॥57\,[58]\!॥
\end{sans}

%58
\asana{Supporting [himself ] with the hands on the ground, placing the knees on the shoulders and the
heels on the chest, [the yogi] should remain thus. [This] is the \engname{Garland Pose} (\engiast{mālāsana}).}
{ಅಂಗೈಗಳನ್ನು ನೆಲಕ್ಕೆ ಊರಿ, ಮೊಣಕಾಲುಗಳನ್ನು ಹೆಗಲ ಮೇಲೂ ಮತ್ತು ಹಿಮ್ಮಡಿಗಳನ್ನು ಎದೆಯ ಮೇಲೂ ಇರಿಸಬೇಕು. ಇದು \tbf{ಹಾರದ ಭಂಗಿ} (\tit{ಮಾಲಾಸನ}).}
\vskip 0.5cm

\insimgcol{HRCP/061_folio_16_verso_58.jpg}

\asananamestmp{mālāsanaṃ}{Garland Pose}{ಮಾಲಾಸನ}
\end{paracol}
\newpage

\insimg[0.77]{HRCP/062_folio_17_recto_59.jpg}

\asananamestmpscale{haṃsāsanaṃ}{Flamingo Pose}{ಹಂಸಾಸನ}{0.77}
\vskip 0.5cm

\begin{adjustwidth}{2.7cm}{2.7cm}
\begin{sans}
kukkuṭāsanavatsthitvā skandhaparyantaṃ ūruṇī\endnote{ūruṇī~] \textit{corr.}~:  ūruṇi M.} nītvā tiṣṭhet\!। haṃsāsanaṃ bhavati\!॥58\,[59]\!॥
\end{sans}

%59
\begin{multicols}{2}
\asana{Positioned as in the Cock Pose and bringing the thighs up to the shoulders, [the yogi] should
remain thus. [This] is the \engname{Flamingo Pose} (\engiast{haṃsāsana}).}
{ಕುಕ್ಕುಟಾಸನದಲ್ಲಿದ್ದಾಗ ತೊಡೆಗಳನ್ನು ಹೆಗಲವರೆಗೆ ಮೇಲೆತ್ತಬೇಕು. ಇದು \tbf{ಬಾತುಕೋಳಿಯ ಭಂಗಿ} (\tit{ಹಂಸಾಸನ}).}
\end{multicols}
\end{adjustwidth}

\newpage

\insimg[0.86]{HRCP/063_folio_17_recto_60.jpg}

\asananamestmpscale{vānarāsanaṃ}{Monkey Pose}{ವಾನರಾಸನ}{0.86}
\vskip 0.5cm

\begin{adjustwidth}{1.7cm}{1.7cm}
\begin{sans}
jānunī bhūmau saṃsthāpya hastābhyāṃ bāhū parasparaṃ dhṛtvā saralaṃ tiṣṭhet\!। vānarāsanaṃ bhavati\!॥59\,[60]\!॥
\end{sans}

%60
\begin{multicols}{2}
\asana{[The yogi] should place the knees on the ground, hold with the hands the arms crossed over one
another and remain upright. [This] is the \engname{Monkey Pose} (\engiast{vānarāsana}).}{ಮೊಣಕಾಲುಗಳನ್ನು ನೆಲಕ್ಕೆ ಊರಿ, ಕೈಗಳನ್ನು ಪರಸ್ಪರ ಅಡ್ಡಲಾಗಿ ಹಿಡಿದು ನೇರವಾಗಿ ನಿಲ್ಲಬೇಕು. ಇದು \tbf{ಮಂಗನ ಭಂಗಿ} (\tit{ವಾನರಾಸನ}).}
\end{multicols}
\end{adjustwidth}

\newpage

\columnratio{0.5, 0.5}
\begin{paracol}{2}
\insimgcol{HRCP/064_folio_17_verso_61.jpg}

\asananamestmp{parvatāsanaṃ}{Mountain Pose}{ಪರ್ವತಾಸನ}
\vskip 0.5cm

\begin{sans}
jaṅghādvayaṃ parasparaṃ veṣṭayitvā bhūmau saṃsthāpya tadupari ūruṇī nidhāya tadupari nitambaṃ nidhāya tiṣṭhet\!। parvatāsanaṃ bhavati\!॥60\,[61]\!॥
\end{sans}

%61
\asana{[The yogi] should wrap the lower legs around each other, place them on the ground, put the thighs
and buttocks on them and remain thus. [This] is the \engname{Mountain Pose} (\engiast{parvatāsana}).}
{ವಾನರಾಸನದಲ್ಲಿದ್ದು ಕಾಲುಗಳ ಕೆಳಭಾಗವನ್ನು ಒಂದರಮೇಲೊಂದಿಟ್ಟು ನೆಲದ ಮೇಲಿರಿಸಿ, ಅದರ ಮೇಲೆ ತೊಡೆ ಮತ್ತು ಪೃಷ್ಠವನ್ನಿಟ್ಟು ಹಾಗೆಯೇ ಕುಳಿತುಕೊಳ್ಳಬೇಕು. ಇದು \tbf{ಪರ್ವತಾಸನ}.}

\switchcolumn
~\vskip 3.2cm

\begin{sans}
pādatale bhūmau nidhāya jānunī urasi nidhāya viparītahastābhyāṃ ūrusahitajaṅghe nibaddhya tiṣṭhet\!। pāśāsanaṃ bhavati\!॥61\,[62]\!॥
\end{sans}

%62
\asana{[The yogi] should place the soles of the feet on the ground and the knees on the chest, and bind
the lower legs and thighs with the hands crossed over. He should remain thus. [This] is the \engname{Noose
Pose} (\engiast{pāśāsana}).}{ಪಾದಗಳನ್ನು ನೆಲಕ್ಕೆ ಊರಿ ಮೊಣಕಾಲುಗಳನ್ನು ಎದೆಗೆ ತಾಗಿಸಬೇಕು. ಬಲಗೈಯನ್ನು ಅಡ್ಡಲಾಗಿ ತಂದು ಕಾಲುಗಳನ್ನು ಬಿಗಿಯಾಗಿ ಅಪ್ಪಿಕೊಳ್ಳಬೇಕು. ಎಡಗೈಯನ್ನು ಬೆನ್ನ ಹಿಂಭಾಗದಿಂದ ಮುಂದಕ್ಕೆ ತಂದು ಬಲಗೈಯೊಂದಿಗೆ ಜೋಡಿಸಬೇಕು. ಇದು \tbf{ಹಗ್ಗದ ಭಂಗಿ} (\tit{ಪಾಶಾಸನ}).}
\vskip 0.5cm

\insimgcol{HRCP/065_folio_17_verso_62.jpg}

\asananamestmp{pāśāsanaṃ}{Noose Pose}{ಪಾಶಾಸನ}
\end{paracol}

\newpage

\columnratio{0.5, 0.5}
\begin{paracol}{2}
\insimgcol{HRCP/066_folio_18_recto_63.jpg}

\asananamestmp{kādambaṃ}{Goose [Pose]}{ಕಾದಂಬಾಸನ}
\vskip 0.5cm

\begin{sans}
pārṣṇibhyāṃ bhūmimavaṣṭabhya hastadvayena gulphau dhṛtvā tiṣṭhet\!। 
kādambaṃ bhavati\!॥62\,[63]\!॥
\end{sans}

%63
\asana{[The yogi] should support himself with the heels on the ground, hold the ankles with the hands and remain thus. [This] is the \engname{Goose [Pose]} (\engiast{kādamba}).}
{ಹಿಮ್ಮಡಿಗಳನ್ನು ನೆಲಕ್ಕೆ ಊರಿ, ಕೈಗಳಿಂದ ಪಾದದ ಗಂಟನ್ನು ಹಿಡಿದು ಕುಳಿತುಕೊಳ್ಳಬೇಕು. ಇದು \tbf{ಕಾದಂಬಾಸನ}.}

\switchcolumn
~\vskip 3.7cm

\begin{sans}
ūrumadhyād hastau niveśya nitambaṃ dhṛtvā pādatalābhyāṃ bhūmimavaṣṭabhya tiṣṭhet\!। kāñcyāsanaṃ bhavati\!॥63\,[64]\!॥
\end{sans}

%64
\asana{Having passed the hands through the thighs, [the yogi] should hold the buttocks, support [himself] with the soles of the feet on the ground and remain thus. [This] is the \engname{Girdle Pose}
(\engiast{kāñcyāsana}).}{ತೊಡೆಗಳ ಮಧ್ಯದಿಂದ ಕೈಗಳನ್ನು ತೂರಿಸಿ ಪೃಷ್ಠದ ಭಾಗವನ್ನು ಹಿಡಿಯಬೇಕು. ಪಾದಗಳನ್ನು ನೆಲಕ್ಕೆ ಊರಿ ಹಾಗೆಯೇ ಇರಬೇಕು. ಇದು  \tbf{ಸೊಂಟದ ಪಟ್ಟಿಯ ಭಂಗಿ} (\tit{ಕಾಂಚ್ಯಾಸನ}).}
\vskip 0.5cm

\insimgcol{HRCP/067_folio_18_recto_64.jpg}

\asananamestmp{kāñcyāsanaṃ}{Girdle Pose}{ಕಾಂಚ್ಯಾಸನ}
\end{paracol}

\newpage

\columnratio{0.5, 0.5}
\begin{paracol}{2}
\insimgcol{HRCP/068_folio_18_verso_65.jpg}

\asananamestmp{aṅgamoṭanaṃ}{Wringing the Limbs [Pose]}{ಅಂಗಮೋಟನಾಸನ}
\vskip 0.5cm

\begin{sans}
hastayoraṅgulīrbaddhvā hastayormadhyātsarvamaṅgaṃ niṣkāsayitvā  tiṣṭhet\!। aṅgamoṭanaṃ\break bhavati\!॥64\,[65]\!॥
\end{sans}

%65
\asana{Having bound the fingers of the hands, [the yogi] should make his whole body pass through the
middle of the arms and remain thus. [This] is the \engname{Wringing the Limbs [Pose]} (\engiast{aṅgamoṭana}).}
{ಕೈಬೆರಳುಗಳನ್ನು ಪರಸ್ಪರ ಕೂಡಿಸಿ, ಇಡೀ ಶರೀರವನ್ನು ಆ ಕೈಗಳ ಮಧ್ಯದಿಂದ ಹಾದುಹೋಗುವಂತೆ ಮಾಡಬೇಕು. ಇದನ್ನು \tbf{ಅಂಗಗಳನ್ನು ಹಿಂಡುವ ಭಂಗಿ} ಎನ್ನಲಾಗುತ್ತದೆ (\tit{ಅಂಗಮೋಟನಾಸನ}).}
\switchcolumn
~\vskip 3.7cm

\begin{sans}
ekaikaṃ pādatalaṃ kukṣau nidhāya tiṣṭhet\!। ucchīrṣakaṃ bhavati\!॥65\,[66]\!॥
\end{sans}

%66
\asana{[The yogi] should place the sole of one foot in the armpit, remain thus, [release it and place the sole
of the other foot in the other armpit]. [This] is the \engname{Pillow [Pose]} (\engiast{ucchīrṣaka}).}
{ಒಂದು ಪಾದದ ಅಡಿಯನ್ನು ಕಂಕುಳಿನಲ್ಲಿ ಇರಿಸಬೇಕು. ನಂತರ ಇದನ್ನೇ ಇನ್ನೊಂದು ಪಾದದಿಂದಲೂ ಮಾಡಬೇಕು. ಇದು \tbf{ದಿಂಬಿನ ಭಂಗಿ} (\tit{ಉಚ್ಚೀರ್ಶಕಾಸನ}).}
\vskip 0.5cm

\insimgcol{HRCP/069_folio_18_verso_66.jpg}

\asananamestmp{ucchīrṣakaṃ}{Pillow [Pose]}{ಉಚ್ಚೀರ್ಶಕಾಸನ}
\end{paracol}

\newpage

\columnratio{0.5, 0.5}
\begin{paracol}{2}
\insimgcol{HRCP/070_folio_19_recto_67.jpg}
%\asananamestmp{}{}{}
\vskip 0.5cm

\begin{sans}
jānuṃ\endnote{jānuṃ~] \textit{corr.}~: jānu M.} stanapradeśe nidhāya pārṣṇiṃ\endnote{pārṣṇiṃ~] \textit{emend.}, pārṣṇī  M.} dvitīyastanapradeśe nidhāya hastadvayena dhṛtvā tiṣṭhet\!॥66\,[67]\!॥
\end{sans}

%67
\asana{[The yogi] should place the knee on [one] side of the chest and the heel [of that leg] on the other
side of the chest. He should hold [the knee and heel] with the hands and remain thus.}
{ಮೊಣಕಾಲನ್ನು ಎದೆಯ ಒಂದು ಬದಿಯಲ್ಲಿ ಮತ್ತು ಹಿಮ್ಮಡಿಯನ್ನು ಎದೆಯ ಇನ್ನೊಂದು ಬದಿಯಲ್ಲಿ ಇರಿಸಿ, ಕೈಗಳಿಂದ ಅವುಗಳನ್ನು ಹಿಡಿದುಕೊಳ್ಳಬೇಕು.}

\switchcolumn
~\vskip 4.7cm

\begin{sans}
hastatalayoḥ pādatale kṛtvā tiṣṭhet\!। gacchet\!। pādukāsanaṃ bhavati\!॥67\,[68]\!॥
\end{sans}

%68
\asana{Having put the soles of the feet on the palms of the hands, [the yogi] should remain thus and
[then] walk around. [This] is the \engname{Sandal Pose} (\engiast{pādukāsana}).}
{ಪಾದಗಳನ್ನು ಅಂಗೈಗಳ ಮೇಲೆ ಇಟ್ಟುಕೊಂಡು ಓಡಾಡಬೇಕು. ಇದು \tbf{ಪಾದರಕ್ಷೆಯ ಭಂಗಿ}\break (\tit{ಪಾದುಕಾಸನ}).}
\vskip 0.5cm

\insimgcol{HRCP/071_folio_19_recto_68.jpg}

\asananamestmp{pādukāsanaṃ}{Sandal Pose}{ಪಾದುಕಾಸನ}
\end{paracol}

\newpage

\columnratio{0.5, 0.5}
\begin{paracol}{2}
\insimgcol{HRCP/072_folio_19_verso_69.jpg}

\asananamestmp{grāhāsanaṃ}{Crocodile Pose}{ಗ್ರಾಹಾಸನ}
\vskip 0.5cm

\begin{sans}
pādatale bhūmau saṃsthāpya jānumadhye kūrparau saṃsthāpya hastābhyāṃ gulphapradeśe dhṛtvā tiṣṭhet\!। grāhāsanaṃ bhavati\!॥68\,[69]\!॥
\end{sans}

%69
\asana{[The yogi] should place the soles of the feet on the ground and the elbows on the insides of the
[bent] knees, hold the region of the ankles with the hands, and remain thus. [This] is the
\engname{Crocodile Pose} (\engiast{grāhāsana}).}
{ಪಾದಗಳನ್ನು ನೆಲಕ್ಕೆ ಊರಿ, ಮಡಚಿದ ಮೊಣಕಾಲುಗಳ ಒಳಗಿನಿಂದ ಮೊಣಕೈಗಳನ್ನು ತಂದು ಪಾದದ ಗಂಟನ್ನು ಹಿಡಿಯಬೇಕು. ಇದು \tbf{ಮೊಸಳೆಯ ಭಂಗಿ} (\tit{ಗ್ರಾಹಾಸನ}).}

\switchcolumn
~\vskip 3.7cm

\begin{sans}
muṣṭī bhūmau saṃsthāpya pādau daṇḍavatprasārya tiṣṭhet\!। parpaṭāsanaṃ bhavati\!॥69\,[70]\!॥
\end{sans}

%70
\asana{Placing the fists on the ground and extending the legs [straight] like a stick, [the yogi] should
remain thus. [This] is the \engname{Parpaṭa-plant Pose} (\engiast{parpaṭāsana}).}{ಮುಷ್ಟಿಗಳನ್ನು ನೆಲಕ್ಕೆ ಊರಿ, ಕೈಗಳ ಮಧ್ಯದಿಂದ ದೇಹವನ್ನು ಮೇಲೆತ್ತಿ, ಕಾಲುಗಳನ್ನು ಕೋಲಿನಂತೆ ನೇರವಾಗಿ ಚಾಚಬೇಕು. ಇದು \tbf{ಪರ್ಪಟಾಸನ}.}
\vskip 0.5cm

\insimgcol{HRCP/073_folio_19_verso_70.jpg}

\asananamestmp{parpaṭāsanaṃ}{Parpaṭa-plant Pose}{ಪರ್ಪಟಾಸನ}
\end{paracol}

\newpage

\insimg[0.83]{HRCP/074_folio_20_recto_71.jpg}

\asananamestmpscale{aśvasādanaṃ gajavadgajasādanaṃ}{Horse’s Seat, Elephant’s Seat}{ಅಶ್ವಸಾಧನಾ, ಗಜಸಾಧನಾ}{0.83}
\vspace{0.5cm}

\begin{adjustwidth}{2cm}{2cm}
\begin{sans}
muṣṭī kaniṣṭhikāpradeśena\endnote{°pradeśena~] M(pc), pradeśe kṛtvā M(ac).} bhūmau nidhāya upari pādatale saṃsthāpya aśvavaccharīraṃ cālanīyaṃ aśvasādanaṃ bhavati\!। gajavadgajasādanaṃ bhavati\!॥70\,[71]\!॥
\end{sans}

%71
\begin{multicols}{2}
\asana{[The yogi] should put the fists with the little fingers on the ground, place the soles of the feet on
[them] and move the body like a horse. [This] is the \engname{Horse’s Seat} (\engiast{aśvasādana}). [Likewise,] the
\engname{Elephant’s Seat} (\engiast{gajasādana}) is [moving the body] like an elephant.}
{ಕಿರುಬೆರಳು ನೆಲಕ್ಕೆ ತಗಲುವಂತೆ ಮುಷ್ಟಿಗಳನ್ನು ಊರಿ, ಅವುಗಳ ಮೇಲೆ ಪಾದಗಳನ್ನಿಟ್ಟು ಕುದುರೆಯಂತೆ ಚಲಿಸಬೇಕು. ಇದು \tbf{ಅಶ್ವಸಾಧನಾ}. ಇದೇ ರೀತಿ ಆನೆಯಂತೆ ಚಲಿಸುವುದು \tbf{ಗಜಸಾಧನಾ}.}
\end{multicols}
\end{adjustwidth}

\newpage

\insimg[0.83]{HRCP/075_folio_20_recto_72.jpg}

\asananamestmpscale{dviśīrṣaṃ}{Two-headed [Pose]}{ದ್ವಿಶೀರ್ಷಾಸನ}{0.83}
\vspace{0.5cm}

\begin{adjustwidth}{2cm}{2cm}
\begin{sans}
skandhaṃ\endnote{skandhaṃ~] \textit{emend.}~: skandhaḥ M.} śiraḥparyantaṃ nītvā sthātavyam\!। dviśīrṣaṃ bhavati\!॥71\,[72]\!॥
\end{sans}

%72
\begin{multicols}{2}
\asana{Having taken [one] shoulder up to the head [while sitting], [the yogi] should remain thus. [This]
is the \engname{Two-headed [Pose]} (\engiast{dviśīrṣa}).}{ಕುಳಿತಿರುವಾಗ ಒಂದು ಹೆಗಲನ್ನು ತಲೆಯವರೆಗೆ ಮೇಲೆತ್ತಿ ಹಾಗೆಯೇ ಇರಬೇಕು. ಇದು \tbf{ದ್ವಿಶೀರ್ಷಾಸನ}.}
\end{multicols}
\end{adjustwidth}

\newpage

\columnratio{0.5, 0.5}
\begin{paracol}{2}
\insimgcol{HRCP/076_folio_20_verso_73.jpg}

\asananamestmp{kubjāsanaṃ}{Humpbacked Pose}{ಕುಬ್ಜಾಸನ}
\vskip 0.5cm

\begin{sans}
nābhau hanuṃ nidhāya tiṣṭhet\!। kubjāsanaṃ bhavati\!॥72\,[73]\!॥
\end{sans}

%73
\asana{[The yogi] should put his jaw on his navel and remain thus. [This] is the \engname{Humpbacked Pose}
(\engiast{kubjāsana}).}{ಗಲ್ಲವನ್ನು ಹೊಕ್ಕುಳಿನ ಮೇಲೆ (ನಾಭಿ) ಇಟ್ಟುಕೊಳ್ಳಬೇಕು. ಇದು \tbf{ಗೂನು ಬೆನ್ನಿನ ಭಂಗಿ} (\tit{ಕುಬ್ಜಾಸನ}).}

\switchcolumn
~\vskip 3.2cm

\begin{sans}
hastatale bhūmau nidhāya hastamadhyāt bhūmiṃ santyajya\endnote{santyajya~] \textit{emend.}~: tyajya M.} daṇḍarūpau pādau kṛtvā preṅkhavaccharīraṃ\endnote{preṅkhavac~] \textit{emend.},  preṅkhava M.} cālanīyam\!। preṅkhāsanaṃ bhavati\!॥73\,[74]\!॥
\end{sans}

%74
\asana{Placing the palms of the hands on the ground, [the yogi] should leave the ground [by lifting the
body up] between the hands, make the legs [straight] like a stick and move the body around like a
swing. [This] is the \engname{Swing Pose} (\engiast{preṅkhāsana}).}
{ಅಂಗೈಗಳನ್ನು ನೆಲಕ್ಕೆ ಊರಿ, ಕೈಗಳ ಮಧ್ಯದ ಜಾಗದಿಂದ ಶರೀರವನ್ನು ಮೇಲೆತ್ತಬೇಕು. ಕಾಲುಗಳನ್ನು ನೇರವಾಗಿ ಚಾಚಿ ತೂಗುವ ಉಯ್ಯಾಲೆಯಂತೆ ಶರೀರವನ್ನು ಆಡಿಸಬೇಕು. ಇದು \tbf{ಉಯ್ಯಾಲೆ ಭಂಗಿ} (\tit{ಪ್ರೇಂಖಾಸನ}).}
\vskip 0.5cm

\insimgcol{HRCP/077_folio_20_verso_74.jpg}

\asananamestmp{preṅkhāsanaṃ}{Swing Pose}{ಪ್ರೇಂಖಾಸನ}
\end{paracol}

\newpage

\insimg[0.75]{HRCP/078_folio_21_recto_75.jpg}
%\asananamestmp{}{}{}
\vspace{0.25cm}

\begin{adjustwidth}{3cm}{3cm}
\begin{sans}
preṅkhavatsthitvā pādāgre mastakopari nītvā pṛṣṭhabhāgaṃ\endnote{pṛṣṭhabhāgaṃ~] \textit{em.}~: pṛṣṭhabhāge M.} bhūmau nidhāya tiṣṭhet\!॥74\,[75]\!॥
\end{sans}

%75
\begin{multicols}{2}
\asana{[The yogi] should remain as in the Swing Pose, take the toes over the head and place the back [of
the body] on the ground, and remain thus.}{ಪ್ರೇಂಖಾಸನದಲ್ಲಿದ್ದಾಗ ಕಾಲ್ಬೆರಳುಗಳನ್ನು ತಲೆಯ ಮೇಲಿಂದ ತಂದು ಬೆನ್ನನ್ನು ನೆಲಕ್ಕೆ ತಾಗಿಸಬೇಕು.}
\end{multicols}
\end{adjustwidth}

\newpage

\insertFullImage{097_folio_25_verso_94_E_STATIONARY.jpg}

\newpage

\pagesWithBorder

\null\vfill
{\leftskip=30pt
\begin{sans}
{\fontsize{18pt}{21.6pt}\selectfont \color{deepmaroon}{atha utthānāsanāni\endnote{utthānāsanāni~] \textit{conj.}~: uttāna āsanāni M.}}}
\end{sans}
\vskip 7pt

\begin{eng}
\prtheading[Standing Poses]{Now, the Standing Poses}
\end{eng}
\vskip 7pt

\begin{kan}
\prtheading[ನಿಂತು ಮಾಡುವ ಆಸನಗಳ ವಿವರ ಇಲ್ಲಿದೆ]{ಈಗ, ನಿಂತು ಮಾಡುವ ಆಸನಗಳ ವಿವರ ಇಲ್ಲಿದೆ}
\end{kan}\par}
\vfill

\newpage


\insimg[0.85]{HRCP/079_folio_21_recto_76.jpg}

\asananamestmpscale{utpīḍāsanaṃ}{Pressure Pose}{ಉತ್ಪೀಡನಾಸನ}{0.85}
\vskip 0.5cm

\begin{adjustwidth}{1.85cm}{1.85cm}
\begin{sans}
pārṣṇidvayaṃ saṃmīlya sthitvā nitambaṃ jānupradeśe ānīya tiṣṭhet\!। utpīḍāsanaṃ bhavati\!॥75\,[76]\!॥
\end{sans}

%76
\begin{multicols}{2}
\asana{Joining the heels together [while] standing, [the yogi] should take the buttocks to the level of the
knees and remain thus. [This] is the \engname{Pressure Pose} (\engiast{utpīḍāsana}).}{ನಿಂತಿರುವ ಸ್ಥಿತಿಯಲ್ಲಿ ಹಿಮ್ಮಡಿಗಳನ್ನು ಪರಸ್ಪರ ಜೋಡಿಸಿ, ಪೃಷ್ಠವನ್ನು ಮೊಣಕಾಲುಗಳ ಮಟ್ಟಕ್ಕೆ ಕೆಳಕ್ಕೆ ತಂದು ಹಾಗೆಯೇ ಇರಬೇಕು. ಇದನ್ನು \tbf{ಒತ್ತಡದ ಭಂಗಿ} ಎನ್ನಲಾಗುತ್ತದೆ (\tit{ಉತ್ಪೀಡಾಸನ}).}
\end{multicols}
\end{adjustwidth}

\newpage

\insimg[0.9]{HRCP/080_folio_21_verso_77.jpg}

\asananamestmpscale{vimānāsanaṃ}{Flying Chariot Pose}{ವಿಮಾನಾಸನ}{0.9}
\vskip 0.5cm

\begin{adjustwidth}{1.2cm}{1.2cm}
\begin{sans}
ekaṃ pādaṃ bhūmau nidhāya nitambaṃ jānupradeśe ānīya dvitīyapādaṃ jānuni nidhāya tiṣṭhet\!। vimānāsanaṃ\endnote{vimānāsanaṃ~] M(ac), vimānasanaṃ P, vimalāsanaṃ M(pc).} bhavati\!॥76\,[77]\!॥
\end{sans}

%77
\begin{multicols}{2}
\asana{Keeping one foot on the ground, [the yogi] should take the buttocks to the level of the knees and
put the other foot on the knee. He should remain thus. [This] is the \engname{Flying Chariot Pose}
(\engiast{vimānāsana}).}{ಒಂದು ಪಾದವನ್ನು ನೆಲದ ಮೇಲಿಟ್ಟು, ಪೃಷ್ಠವನ್ನು ಮೊಣಕಾಲುಗಳ ಮಟ್ಟಕ್ಕೆ ತರಬೇಕು. ಇನ್ನೊಂದು ಪಾದವನ್ನು ಆ ಮೊಣಕಾಲಿನ ಮೇಲಿಟ್ಟು ಹಾಗೆಯೇ ಇರಬೇಕು. ಇದು \tbf{ವಿಮಾನಾಸನ}.}
\end{multicols}
\end{adjustwidth}

\newpage

\insimg[0.75]{HRCP/081_folio_21_verso_78.jpg}

\asananamestmpscale{kapotanāṭakaṃ}{[Pose of] the Pigeondancer}{ಕಪೋತನಾಟಕಾಸನ}{0.75}
\vskip 0.5cm

\begin{adjustwidth}{3cm}{3cm}
\begin{sans}
pādatale bhūmau nidhāya pṛṣṭhabhāgena hastau jaṅghāparyantaṃ nītvā  tiṣṭhet\!। kapotanāṭakaṃ bhavati\!॥77\,[78]\!॥
\end{sans}

%78
\begin{multicols}{2}
\asana{Keeping the soles of the feet on the ground, [the yogi] should bring the hands [down] along the
back [of the body] as far as the lower legs and remain thus. [This] is the \engname{[Pose of ] the Pigeondancer}
(\engiast{kapotanāṭaka}).}{ಪಾದಗಳನ್ನು ನೆಲಕ್ಕೆ ಊರಿ ನಿಂತು, ಕೈಗಳನ್ನು ಬೆನ್ನಿನ ಹಿಂದಿನಿಂದ ಕಾಲುಗಳ ಕೆಳಭಾಗದವರೆಗೆ ತಂದು ಹಾಗೆಯೇ ಇರಬೇಕು. ಇದು \tbf{ಪಾರಿವಾಳದ ನೃತ್ಯದ ಭಂಗಿ} (\tit{ಕಪೋತನಾಟಕಾಸನ}).}
\end{multicols}
\end{adjustwidth}

\newpage

\insimg[0.75]{HRCP/082_folio_22_recto_79.jpg}

\asananamestmpscale{ardhacandraṃ}{Half-moon [Pose]}{ಅರ್ಧಚಂದ್ರಾಸನ}{0.75}
\vskip 0.5cm

\begin{adjustwidth}{2.95cm}{2.95cm}
\begin{sans}
ekena pādena bhūmimavaṣṭabhya anyapādatalena ūrumavaṣṭabhya\endnote{ūrum~] \textit{emend}., ūru M.} tiṣṭhedāsta\endnote{āsta~] \textit{conj.}, āṣṭa  P, āniṣṭha M.} iti punaḥ punaḥ kartavyam\!। ardhacandraṃ bhavati\!॥78\,[79]\!॥
\end{sans}

%79
\begin{multicols}{2}
\asana{Supporting [himself ] with one foot on the ground and with the sole of the other foot on the thigh,
[the yogi] should stand [and] sit. It should be done thus again and again. [This] is the \engname{Half-moon [Pose]} (\engiast{ardhacandra}).}{ಒಂದು ಪಾದವನ್ನು ನೆಲದ ಮೇಲಿಟ್ಟು, ಇನ್ನೊಂದು ಪಾದದ ಅಡಿಯನ್ನು ನೆಲದ ಮೇಲಿಟ್ಟ ಕಾಲಿನ ತೊಡೆಯ ಮೇಲಿರಿಸಬೇಕು. ಈ ಸ್ಥಿತಿಯಲ್ಲೇ ಕುಳಿತುಕೊಳ್ಳುವುದು ಮತ್ತು ಏಳುವುದು—ಹೀಗೆ ಪದೇ ಪದೇ ಮಾಡಬೇಕು. ಇದು \tbf{ಅರ್ಧಚಂದ್ರಾಸನ}.}
\end{multicols}
\end{adjustwidth}

\newpage

\insimg[0.8]{HRCP/083_folio_22_recto_80.jpg}

\asananamestmpscale{śaṅkvāsanaṃ}{Spike Pose}{ಶಂಕ್ವಾಸನ}{0.8}
\vskip 0.5cm

\begin{adjustwidth}{2.4cm}{2.4cm}
\begin{sans}
tiṣṭhansan pārṣṇinā itarajaghanapradeśe saṃpīḍya sthātavyamutthitavyamiti punaḥ punaḥ\endnote{punaḥ~] M(pc), \textit{omitted} M(ac).} kartavyam\!। śaṅkvāsanaṃ bhavati\!॥79\,[80]\!॥
\end{sans}

%80
\begin{multicols}{2}
\asana{While standing, [the yogi] should press the heel [of one foot] on the loin of the other [leg], sit
down and stand up. He should do thus again and again. [This] is the \engname{Spike Pose} (\engiast{śaṅkvāsana}).}
{ನಿಂತಿರುವಾಗ ಒಂದು ಕಾಲಿನ ಹಿಮ್ಮಡಿಯಿಂದ ಇನ್ನೊಂದು ಕಾಲಿನ ಒಳತೊಡೆಯ ಭಾಗವನ್ನು ಒತ್ತಿ ಹಿಡಿದು, ಕುಳಿತು ಏಳಬೇಕು. ಇದನ್ನು ಪದೇ ಪದೇ ಮಾಡಬೇಕು. ಇದು \tbf{ಶಂಕ್ವಾಸನ}.}
\end{multicols}
\end{adjustwidth}

\newpage

{\setlength{\columnsep}{0.7cm}
\columnratio{0.8, 0.2}
\begin{paracol}{2}
~\vskip 1.25cm
\insimgcol{HRCP/084_folio_22_verso_81.jpg}

\asananamestmp{tāṇḍavāsanaṃ}{Pose of [Śiva’s] Tāṇḍava Dance}{ಶಿವನ ತಾಂಡವಾಸನ}
\switchcolumn
~\vfill
\begin{sans}
ekena pādena sthātavyamutthātavyam\!। tāṇḍavāsanaṃ bhavati\!॥80\,[81]\!॥
\end{sans}
\bigskip

%81
\asana{[The yogi] should sit and stand up on one leg. [This] is the \engname{Pose of [Śiva’s]\break Tāṇḍava Dance}\\
(\engiast{tāṇḍavāsana}).}{ಕೇವಲ ಒಂದು ಕಾಲಿನ ಮೇಲೆ ಕುಳಿತು ಏಳಬೇಕು. ಇದನ್ನು \tbf{ಶಿವನ ತಾಂಡವಾಸನ} ಎನ್ನಲಾಗುತ್ತದೆ.}
\vfill
\end{paracol}}

\newpage

{\setlength{\columnsep}{0.7cm}
\columnratio{0.18, 0.82}
\begin{paracol}{2}
~\vskip 5cm
\begin{sans}
grīvāyāṃ pādaṃ saṃsthāpya sthātavyamutthātavyam\!। trivikramāsanaṃ\\ bhavati\!॥81\,[82]\!॥
\end{sans}

%82
\asana{Placing one foot on the [back of the] neck, [the yogi] should sit and stand up. [This] is the \engname{Pose of
[Viṣṇu’s] Three Steps}\\ (\engiast{trivikramāsana}).}{ಒಂದು ಪಾದವನ್ನು ಕುತ್ತಿಗೆಯ\break ಹಿಂಭಾಗದಲ್ಲಿ ಇರಿಸಿಕೊಂಡು,\break ಕುಳಿತು ಏಳಬೇಕು. ಇದು \tbf{ವಿಷ್ಣುವಿನ ಮೂರು ಹೆಜ್ಜೆಗಳ ಭಂಗಿ}\break (\tit{ತ್ರಿವಿಕ್ರಮಾಸನ}).}

\switchcolumn
~\vskip 1.7cm

\insimgcol{HRCP/085_folio_22_verso_82.jpg}

\asananamestmp{trivikramāsanaṃ}{Pose of [Viṣṇu’s] Three Steps}{ತ್ರಿವಿಕ್ರಮಾಸನ}
\end{paracol}}
\newpage

{\setlength{\columnsep}{0.7cm}
\columnratio{0.8, 0.2}
\begin{paracol}{2}
~\vskip 0.75cm
\insimgcol{HRCP/086_folio_23_recto_83.jpg}

\asananamestmp{utthānotthānaṃ}{Standing Up Repeatedly [Pose]}{ಉತ್ಥಾನೋತ್ತಾನಾಸನ}

\switchcolumn

~\vfill
\begin{sans}
punaḥ punarutthātavyaṃ āsitavyam\!। utthānotthānaṃ bhavati\!॥82\,[83]\!॥
\end{sans}

%83
\asana{\raggedright Again and again, [the yogi] should stand up and sit down. [This] is the \engname{Standing Up Repeatedly
[Pose]} (\engiast{utthānotthāna}).}{\raggedright ಯಾವುದೇ ಆಧಾರವಿಲ್ಲದೆ ಪದೇ ಪದೇ ಕುಳಿತು ಏಳುತ್ತಿರಬೇಕು. ಇದನ್ನು \tbf{ಉತ್ಥಾನೋತ್ತಾನಾಸನ} ಎನ್ನಲಾಗುತ್ತದೆ.}
\vfill
\end{paracol}}

\newpage

\insimg[0.83]{HRCP/087_folio_23_recto_84.jpg}

\asananamestmpscale{āliṅganāsanaṃ}{Embracing [the Wall] Pose}{ಆಲಿಂಗನಾಸನ}{0.83}
\vskip 0.5cm

\begin{adjustwidth}{2cm}{2cm}
\begin{sans}
hastatrayaṃ bhittiṃ parityajya sthātavyam। bhittau hṛdayaṃ saṃspṛśya niṣkāsya punaḥ saṃspṛśet। āliṅganāsanaṃ\endnote{āliṅganāsanaṃ~] \textit{emend.}, āliṅgāsanaṃ M.} bhavati॥\!83 [84]\!॥
\end{sans}

%84
\begin{multicols}{2}
\asana{Standing at a distance of three cubits from a wall, [the yogi] should touch his chest on the wall,
then push it away. He should touch it thus again [and again]. [This] is the \engname{Embracing [the Wall] Pose} (\engiast{āliṅganāsana}).}{ಗೋಡೆಯಿಂದ ಮೂರು ಗೇಣುಗಳಷ್ಟು ದೂರದಲ್ಲಿ ನಿಂತು, ಎದೆಯನ್ನು ಗೋಡೆಗೆ ತಾಗಿಸಬೇಕುನಂತರ ಹಿಂದಕ್ಕೆ ತಳ್ಳಬೇಕು. ಇದನ್ನು ಪದೇ ಪದೇ ಮಾಡಬೇಕು. ಇದು  \tbf{ಗೋಡೆಯನ್ನು ಅಪ್ಪಿಕೊಳ್ಳುವ ಭಂಗಿ} (\tit{ಆಲಿಂಗನಾಸನ}).}
\end{multicols}
\end{adjustwidth}

\newpage

{\setlength{\columnsep}{0.7cm}
\columnratio{0.8, 0.2}
\begin{paracol}{2}
~\vskip 0.5cm
\insimgcol{HRCP/088_folio_23_verso_85.jpg}

\asananamestmp{bālāliṅganaṃ}{Embracing a Child [Pose]}{ಬಾಲಾಲಿಂಗನಾಸನ}

\switchcolumn
~\vfill
\begin{sans}
ekaṃ jānum\endnote{jānum~] \textit{emend.}~: jānu M.} urasi āliṅgya sthātavyam\!। bālāliṅganaṃ bhavati\!॥84\,[85]\!॥
\end{sans}

%85
\asana{\raggedright Hugging one knee to the chest, [the yogi] should stand. [This] is the \engname{Embracing a Child [Pose]}
(\engiast{bālāliṅgana}).}{\raggedright  ಒಂದು ಮೊಣಕಾಲನ್ನು ಎದೆಗೆ ಅಪ್ಪಿಕೊಂಡು ನಿಲ್ಲಬೇಕು. ಇದು \tbf{ಮಗುವನ್ನು ಅಪ್ಪಿಕೊಳ್ಳುವ ಭಂಗಿ} (\tit{ಬಾಲಾಲಿಂಗನಾಸನ}).}
\vfill
\end{paracol}}

\newpage

{\setlength{\columnsep}{0.7cm}
\columnratio{0.2, 0.8}
\begin{paracol}{2}
~\vskip 3.7cm
\begin{sans}
vṛṣaṇasahitaṃ liṅgaṃ ūrumadhye gāḍhaṃ dhṛtvā caraṇāgrābhyāṃ sthātavyam\!। kaupīnāsanaṃ bhavati\!॥85\,[86]\!॥
\end{sans}

%86
\asana{\raggedright Holding the penis and scrotum firmly between the thighs, [the yogi] should stand on the tips of
the toes. [This] is the \engname{Loincloth Pose} (\engiast{kaupīnāsana}).}{\raggedright  ಜನನಾಂಗ ಮತ್ತು ಅಂಡಕೋಶಗಳನ್ನು ತೊಡೆಗಳ ಮಧ್ಯೆ ಗಟ್ಟಿಯಾಗಿ ಒತ್ತಿ ಹಿಡಿದು, ಕಾಲ್ಬೆರಳುಗಳ ತುದಿಯ ಮೇಲೆ ನಿಲ್ಲಬೇಕು. ಇದು \tbf{ಕೌಪೀನಾಸನ}.}

\switchcolumn
~\vskip 2cm
\insimgcol{HRCP/089_folio_23_verso_86.jpg}

\asananamestmp{kaupīnāsanaṃ}{Loincloth Pose}{ಕೌಪೀನಾಸನ}

\end{paracol}}


\newpage

\columnratio{0.5, 0.5}
\begin{paracol}{2}

\insimgcol{HRCP/090_folio_24_recto_87.jpg}

\asananamestmp{dehalyullaṅghanaṃ}{Jumping over the Threshold [Pose]}{ದೇಹಲ್ಯುಲ್ಲಂಘನಾಸನ}
\vskip 0.5cm

\begin{sans}
hastadvayaṃ baddhvā tanmadhyāccaraṇadvayaṃ uḍḍānena bahirānīya antarnayet\!।\endnote{antar~] \textit{corr.}~: antaḥ M.} dehalyullaṅghanaṃ bhavati\!॥86\,[87]\!॥
\end{sans}

%87
\asana{Clasping the hands together, [the yogi] should jump the feet through them, outside and inside.
[This] is the \engname{Jumping over the Threshold [Pose]} (\engiast{dehalyullaṅghana}).}{ಕೈಗಳನ್ನು ಪರಸ್ಪರ ಕೂಡಿಸಿ ಹಿಡಿದು, ಪಾದಗಳನ್ನು ಆ ಕೈಗಳ ಒಳಗೆ ಮತ್ತು ಹೊರಗೆ ಹಾರಿಸಬೇಕು. ಇದು \tbf{ಹೊಸ್ತಿಲು ದಾಟುವ ಭಂಗಿ} (\tit{ದೇಹಲ್ಯುಲ್ಲಂಘನಾಸನ}).}

\switchcolumn
~\vskip 4.5cm
\begin{sans}
uḍḍānaṃ kṛtvā pārṣṇibhyāṃ nitambaṃ tāḍayet\!। hariṇāsanaṃ bhavati\!॥87\,[88]\!॥
\end{sans}

%88
\asana{Jumping up, [the yogi] should strike his buttocks with the heels. [This] is the \engname{Deer Pose}
(\engiast{hariṇāsana}).}{ಮೇಲಕ್ಕೆ ಜಿಗಿಯುತ್ತಾ ಹಿಮ್ಮಡಿಗಳಿಂದ ಪೃಷ್ಠಕ್ಕೆ ಹೊಡೆದುಕೊಳ್ಳಬೇಕು. ಇದು \tbf{ಜಿಂಕೆಯ ಭಂಗಿ} (\tit{ಹರಿಣಾಸನ}).}
\vskip 0.5cm

\insimgcol{HRCP/091_folio_24_recto_88.jpg}

\asananamestmp{hariṇāsanaṃ}{Deer Pose}{ಹರಿಣಾಸನ}

\end{paracol}

\newpage

\columnratio{0.5, 0.5}
\begin{paracol}{2}
\insimgcol{HRCP/092_folio_24_verso_89.jpg}

\asananamestmp{musalāsanaṃ}{Pestle Pose}{ಮುಸಲಾಸನ}
\vskip 0.5cm

\begin{sans}
saralaṃ sthitvā ūrdhvaṃ bāhū nayet\!। vāraṃ vāramuḍḍānaṃ kṛtvānayet\!। musalāsanaṃ\break bhavati\!॥88\,[89]\!॥
\end{sans}


%89
\asana{Standing straight, [the yogi] should hold up the arms. Jumping up again and again, he should bring
them down. [This] is the \engname{Pestle Pose}\break (\engiast{musalāsana}).}{ನೇರವಾಗಿ ನಿಂತು ಕೈಗಳನ್ನು ಮೇಲಕ್ಕೆತ್ತಿ ಪದೇ ಪದೇ ಜಿಗಿಯುತ್ತಾ ಕೈಗಳನ್ನು ಕೆಳಕ್ಕೆ ತರಬೇಕು. ಇದು \tbf{ಒನಕೆಯ ಭಂಗಿ} (\tit{ಮುಸಲಾಸನ}).}

\switchcolumn
~\vskip 3.7cm

\begin{sans}
ekena hastena daṇḍarūpasya ekasya pādasyāgraṃ dhṛtvā itarapādatalaṃ bhūmau nidhāya 
tvarayā bhrāmaṇaṃ kartavyam\!। dhruvāsanaṃ bhavati\!॥89\,[90]\!॥
\end{sans}
\bigskip

%90
\asana{Holding with one hand the toes of one leg, which is [straight] like a stick, and placing the sole of
the other foot on the ground, [the yogi] should spin around quickly. [This] is the \engname{Pole Star Pose}
(\engiast{dhruvāsana}).}{ಬಲ ಪಾದವನ್ನು ನೆಲದ ಮೇಲಿಟ್ಟು, ಎಡಗಾಲಿನ ಬೆರಳುಗಳನ್ನು ಎಡಗೈಯಿಂದ ಹಿಡಿದು, ಎಡಗಾಲನ್ನು ಕೋಲಿನಂತೆ ನೇರವಾಗಿಸಿ ವೇಗವಾಗಿ ಸುತ್ತಬೇಕು. ಇದು \tbf{ಧ್ರುವಾಸನ}.}
\vskip 0.5cm

\insimgcol{HRCP/093_folio_24_verso_90.jpg}

\asananamestmp{dhruvāsanaṃ}{Pole Star Pose}{ಧ್ರುವಾಸನ}
\end{paracol}

\newpage

\insimg[0.79]{HRCP/094_folio_25_recto_91.jpg}

\asananamestmpscale{kulālacakrāsanaṃ}{Potter’s Disk Pose}{ಕುಲಾಲ ಚಕ್ರಾಸನ}{0.79}
\vskip 0.5cm

\begin{adjustwidth}{2.5cm}{2.5cm}
\begin{sans}
hastau prasārya bhrāmaṇaṃ kartavyam\!। kulālacakrāsanaṃ bhavati\!॥90\,[91]\!॥
\end{sans}

%91
\begin{multicols}{2}
\asana{Having extended the hands [out to the sides], [the yogi] should spin around. [This] is the \engname{Potter’s
Disk Pose} (\engiast{kulālacakrāsana}).}{ಕೈಗಳನ್ನು ಪಕ್ಕಕ್ಕೆ ಚಾಚಿ ವೇಗವಾಗಿ ಗಿರಗಿರನೆ ಸುತ್ತಬೇಕು. ಇದು \tbf{ಕುಂಬಾರನ ಚಕ್ರದ ಭಂಗಿ} (\tit{ಕುಲಾಲ ಚಕ್ರಾಸನ}).}
\end{multicols}
\end{adjustwidth}
\newpage

\insimg[0.75]{HRCP/095_folio_25_recto_92.jpg}

\asananamestmpscale{uṣṭrāsanaṃ}{Camel Pose}{ಉಷ್ಟ್ರಾಸನ}{0.75}
\vskip 0.5cm

\begin{adjustwidth}{3cm}{3cm}
\begin{sans}
pādāṅguṣṭhābhyāṃ bhūmiṃ dhṛtvā bāhū ūrdhvīkṛtya tiṣṭhet\!। uṣṭrāsanaṃ bhavati\!॥91\,[92]\!॥
\end{sans}

%92
\begin{multicols}{2}
\asana{Having put the big toes of the feet on the ground, [the yogi] should raise up the arms and remain
thus. [This] is the \engname{Camel Pose} (\engiast{uṣṭrāsana}).}{ಕಾಲ್ಬೆರಳುಗಳ ಮೇಲೆ ನಿಂತು, ಕೈಗಳನ್ನು ಮೇಲಕ್ಕೆತ್ತಿ ಹಾಗೆಯೇ ಇರಬೇಕು. ಇದು \tbf{ಒಂಟೆಯ ಭಂಗಿ} (\tit{ಉಷ್ಟ್ರಾಸನ}).}
\end{multicols}
\end{adjustwidth}
\newpage

\insimg[0.8]{HRCP/096_folio_25_verso_93.jpg}

\asananamestmpscale{ākāśakapotaṃ}{Pigeon in the Sky [Pose]}{ಆಕಾಶಕಪೋತಾಸನ}{0.8}
\vskip 0.5cm

\begin{adjustwidth}{2.4cm}{2.4cm}
\begin{sans}
uṣṭrāsane sthitvā caraṇau bhūmerutthāpya mastakopari nītvā pṛṣṭhadeśena bhūmau sthāpayet\!। ākāśakapotaṃ\endnote{ākāśa°~] P, dākāśa° M.} bhavati\!॥92\,[93]\!॥
\end{sans}

%93
\begin{multicols}{2}
\asana{Remaining in Camel Pose, raising the feet from the ground and taking [them] over the head, [the
yogi] should place his back on the ground. This is the \engname{Pigeon in the Sky [Pose]} (\engiast{ākāśakapota}).}
{ಉಷ್ಟ್ರಾಸನದಲ್ಲಿದ್ದಾಗ ಪಾದಗಳನ್ನು ನೆಲದಿಂದ ಮೇಲೆತ್ತಿ ತಲೆಯ ಮೇಲಿಂದ ತಂದು ಬೆನ್ನನ್ನು ನೆಲಕ್ಕೆ ತಾಗಿಸಬೇಕು. ಇದು ಆಕಾಶದ \tbf{ಪಾರಿವಾಳದ ಭಂಗಿ} (\tit{ಆಕಾಶಕಪೋತಾಸನ}).}
\end{multicols}
\end{adjustwidth}
\newpage

\insimg[0.98]{HRCP/097_folio_25_verso_94_E_STATIONARY.jpg}

\asananamestmpscale{garuḍāsanaṃ}{Garuḍa’s Pose}{ಗರುಡಾಸನ}{0.98}
\vskip 0.5cm

\begin{adjustwidth}{0.3cm}{0.3cm}
\begin{sans}
ūrumūle itarapādasya kaniṣṭhikāpradeśe gulphaṃ saṃsthāpya tadeva jānum\endnote{jānum~] \textit{emend.}, jānu M.} itarapādapārṣṇyāṃ\endnote{itarapādapārṣṇyāṃ~] P, itarapārṣṇyāṃ M.} saṃsthāpya tiṣṭhet\!। hastau ca saṃmīlayet\!। garuḍāsanaṃ bhavati\!॥93\,[94]\!॥
\end{sans}

%94
\begin{multicols}{2}
\asana{Placing [the part of ] the ankle on the little toe side of one foot on the [other] thigh, and placing
the knee [of that leg] on the heel of the other foot, [the yogi] should remain thus and join his
hands together. [This] is \engname{Garuḍa’s Pose} (\engiast{garuḍāsana}).}{ಒಂದು ಪಾದದ ಪಾರ್ಶ್ವದ ಭಾಗವನ್ನು ಇನ್ನೊಂದು ಕಾಲಿನ ತೊಡೆಯ ಮೇಲಿಟ್ಟು, ಆ ಕಾಲಿನ ಮೊಣಕಾಲನ್ನು ನಿಂತಿರುವ ಕಾಲಿನ ಹಿಮ್ಮಡಿಯ ಮೇಲಿರಿಸಿ ಕೈಗಳನ್ನು ಜೋಡಿಸಬೇಕು. ಇದು \tbf{ಗರುಡಾಸನ}.}
\end{multicols}
\end{adjustwidth}

\newpage

\insertFullImage{107_folio_28_recto_104_E_ROPES.jpg}

\newpage

\pagesWithBorder

\null\vfill
{\leftskip=30pt
\begin{sans}
{\fontsize{18pt}{21.6pt}\selectfont \color{deepmaroon}{atha rajjvāsanāni\endnote{rajjvāsanāni~] \textit{corr.}, rajvāsanāni M.}}}
\end{sans}
\vskip 7pt

\begin{eng}
\prtheading[Poses with a Rope]{Now, the Poses with a Rope}
\end{eng}
\vskip 7pt

\begin{kan}
\prtheading[ಹಗ್ಗದ ನೆರವಿನಿಂದ ಮಾಡುವ ಆಸನಗಳ ವಿವರಣೆ ಇಲ್ಲಿದೆ]{ಈಗ, ಹಗ್ಗದ ನೆರವಿನಿಂದ ಮಾಡುವ ಆಸನಗಳ ವಿವರಣೆ ಇಲ್ಲಿದೆ}
\end{kan}\par}
\vfill

\newpage

\insimg[1]{HR_not_cropped/098_folio_26_recto_95_S_ROPES.jpg}

\asananamestmp{paroṣṇyāsanaṃ}{Cockroach Pose}{ಪರೋಶ್ಣ್ಯಾಸನ (ಜಿರಳೆಯ ಭಂಗಿ)}
\vskip 1cm

\begin{sans}
% new paragraph
hastadvayena rajjuṃ dhṛtvā pādadvayaṃ hastamadhyātśirasopari bhūmau nidhāya punaḥ punaḥ saṅkṣipet\!। paroṣṇyāsanaṃ bhavati\endnote{bhavati~] P, \textit{omitted} M.}\!॥94\,[95]\!॥
\end{sans}
\bigskip

%95
\asana{Clasping a rope [secured horizontally above the head] with the hands, [the yogi] should pass the
legs between the hands and above the head, and [put them] on the ground. He should propel [the
legs thus] again and again. [This] is the \engname{Cockroach Pose} (\engiast{paroṣṇyāsana}).}{ತಲೆಯ ಮೇಲ್ಭಾಗದಲ್ಲಿ ಅಡ್ಡಲಾಗಿ ಕಟ್ಟಿರುವ ಹಗ್ಗವನ್ನು ಕೈಗಳಿಂದ ಹಿಡಿದುಕೊಂಡು, ಕೈಗಳ ಮಧ್ಯದ ಜಾಗದಿಂದ ಕಾಲುಗಳನ್ನು ತೂರಿಸಿ, ತಲೆಯ ಮೇಲಿಂದ ತಂದು ನೆಲಕ್ಕೆ ಮುಟ್ಟಿಸಬೇಕು. ಇದನ್ನು ಪದೇ ಪದೇ ಮಾಡಬೇಕು. ಇದು \tbf{ಜಿರಳೆಯ ಭಂಗಿ}. (\tit{ಪರೋಶ್ಣ್ಯಾಸನ})}

\newpage

\insimg[0.7]{HR_not_cropped/099_folio_26_recto_96.jpg}
\asananamestmp{daṇḍāsanaṃ}{Stick Pose}{ದಂಡಾಸನ (ಕೋಲಿನ ಭಂಗಿ)}

\begin{sans}
nābhipradeśaṃ\endnote{°deśaṃ~] P,  °deśe M.} rajjau dhārayitvā daṇḍavatsthirībhavet\endnote{sthirī°~] \textit{emend}, sthiraṃ M.}। daṇḍāsanaṃ bhavati\endnote{daṇḍāsanaṃ bhavati~] \textit{diagnostic conj.}, \textit{omitted} M.}\!॥95\,[96]\!॥
\end{sans}
%\bigskip

%96
\asana{Supporting the navel area on a [horizontal] rope, [the yogi] should become [straight and] rigid like
a stick. [This] is the \engname{Stick Pose} (\engiast{daṇḍāsana}).}{ಅಡ್ಡಲಾಗಿರುವ ಹಗ್ಗದ ಮೇಲೆ ಹೊಕ್ಕುಳಿನ ಭಾಗವನ್ನು ಇರಿಸಿ, ಇಡೀ ಶರೀರವನ್ನು ಕೋಲಿನಂತೆ ನೇರವಾಗಿ ಮತ್ತು ಬಿಗಿಯಾಗಿ ಹಿಡಿಯಬೇಕು. ಇದು \tbf{ಕೋಲಿನ ಭಂಗಿ}. (\tit{ದಂಡಾಸನ})}

%\newpage

\insimg[0.7]{HR_not_cropped/100_folio_26_verso_97.jpg}
\asananamestmp{bhārāsanaṃ}{Weight Pose}{ಭಾರಾಸನ (ತೂಕದ ಭಂಗಿ)}

\begin{sans}
rajjau nitambaṃ sthāpayitvā daṇḍavatsthirībhavet\!। bhārāsanaṃ bhavati\!॥96\,[97]\!॥
\end{sans}
%\bigskip

%97
\asana{Having placed the buttocks on a [horizontal] rope, [the yogi] should become [straight and] rigid
like a stick. [This] is the \engname{Weight Pose} (\engiast{bhārāsana}).}{ಅಡ್ಡವಾಗಿರುವ ಹಗ್ಗದ ಮೇಲೆ ಪೃಷ್ಠದ ಭಾಗವನ್ನು ಇರಿಸಿ, ಇಡೀ ಶರೀರವನ್ನು ಕೋಲಿನಂತೆ ನೇರವಾಗಿಸಬೇಕು. ಇದು \tbf{ತೂಕದ ಭಂಗಿ}. (\tit{ಭಾರಾಸನ})}

\newpage

\insimg[0.5]{HR_not_cropped/101_folio_26_verso_98.jpg}
\asananamestmp{nāradāsanaṃ}{Nārada’s Pose}{ನಾರದಾಸನ}

\begin{sans}
hastābhyāṃ rajjuṃ dhṛtvā ūrdhvamārohet\!। nāradāsanaṃ bhavati\!॥97\,[98]\!॥
\end{sans}
%\bigskip

%98
\asana{Holding a [vertical] rope with the hands, [the yogi] should climb up it. [This] is \engname{Nārada’s Pose}
(\engiast{nāradāsana}).}{ನೇರವಾಗಿ ಇಳಿಬಿಟ್ಟಿರುವ ಹಗ್ಗವನ್ನು ಕೈಗಳಿಂದ ಹಿಡಿದು ಮೇಲಕ್ಕೆ ಏರಬೇಕು. ಇದು \tbf{ನಾರದಾಸನ}.}

%\newpage

\insimg[0.5]{HR_not_cropped/102_folio_27_recto_99.jpg}
\asananamestmp{svargāsanaṃ}{[Climbing up to] Heaven Pose}{ಸ್ವರ್ಗಾಸನ}

\begin{sans}
padmāsanaṃ kṛtvā hastābhyāṃ rajjuṃ dhṛtvā  ārohet\!। svargāsanaṃ bhavati\!॥98\,[99]\!॥
\end{sans}
%\bigskip

%99
\asana{Assuming the Lotus Pose and holding a [vertical] rope with the hands, [the yogi] should climb up
it. [This] is the \engname{[Climbing up to] Heaven Pose} (\engiast{svargāsana}).}{ಪದ್ಮಾಸನ ಹಾಕಿದ ಸ್ಥಿತಿಯಲ್ಲೇ ಇಳಿಬಿಟ್ಟಿರುವ ಹಗ್ಗವನ್ನು ಕೈಗಳಿಂದ ಹಿಡಿದು ಮೇಲಕ್ಕೆ ಏರಬೇಕು. ಇದು \tbf{ಸ್ವರ್ಗಾಸನ}.}

\newpage

\long\def\asana#1#2{%
\incanum%
\begin{multicols}{2}
\begin{eng}
\begin{enumerate}[leftmargin=*,labelindent=0pt,widest=999]
\item[\asananum]
#1
\end{enumerate}
\end{eng}

\begin{kan}
\begin{enumerate}[leftmargin=*,labelindent=0pt,widest=999]
\item[\asananum]
#2
\end{enumerate}
\end{kan}
\end{multicols}
}

\insimg[0.5]{HR_not_cropped/103_folio_27_recto_100.jpg}
\asananamestmp{ūrṇanābhyāsanaṃ}{Spider Pose}{ಊರ್ಣನಾಭ್ಯಾಸನ (ಜೇಡದ ಭಂಗಿ)}
%\vskip 1cm

\begin{sans}
kukkuṭāsanaṃ kṛtvā hastābhyāṃ rajjuṃ dhṛtvā ārohet\!। ūrṇanābhyāsanaṃ bhavati\!॥99\,[100]\!॥
\end{sans}
%\bigskip

%100
\asana{Assuming the Cock Pose and holding a [vertical] rope with the hands, [the yogi] should climb up
it. [This] is the \engname{Spider Pose} (\engiast{ūrṇanābhyāsana}).}{ಕುಕ್ಕುಟಾಸನದಲ್ಲಿದ್ದು ಇಳಿಬಿಟ್ಟಿರುವ ಹಗ್ಗವನ್ನು ಕೈಗಳಿಂದ ಹಿಡಿದು ಮೇಲಕ್ಕೆ ಏರಬೇಕು. ಇದು \tbf{ಜೇಡದ ಭಂಗಿ}. (\tit{ಊರ್ಣನಾಭ್ಯಾಸನ})}

%\newpage

\insimg[0.5]{HR_not_cropped/104_folio_27_verso_101.jpg}
\asananamestmp{śukāsanaṃ}{Parrot Pose}{ಶುಕಾಸನ (ಗಿಳಿಯ ಭಂಗಿ)}
%\vskip 1cm

\begin{sans}
muṣṭibhyāṃ rajjuṃ dhṛtvā tadupari pādatale sthāpayitvā tiṣṭhet\!। śukāsanaṃ bhavati\!॥100\,[101]\!॥
\end{sans}
%\bigskip

%101
\asana{Holding a [vertical] rope with the fists, [the yogi] should put the soles of the feet on the [fists]. He
should remain thus. [This] is the \engname{Parrot Pose} (\engiast{śukāsana}).}{ಇಳಿಬಿಟ್ಟಿರುವ ಹಗ್ಗವನ್ನು ಮುಷ್ಟಿಯಿಂದ ಹಿಡಿದು, ಆ ಮುಷ್ಟಿಯ ಮೇಲೆ ಪಾದಗಳನ್ನು ಇರಿಸಿ ಅದೇ ಭಂಗಿಯಲ್ಲಿರಬೇಕು. ಇದು \tbf{ಗಿಳಿಯ ಭಂಗಿ}. (\tit{ಶುಕಾಸನ})}

\newpage

\insimg[0.5]{HR_not_cropped/105_folio_27_verso_102.jpg}
\asananamestmp{tṛṇajalūkaṃ}{Caterpillar [Pose]}{ತೃಣಜಲೂಕಾಸನ (ಕಂಬಳಿ ಹುಳುವಿನ ಭಂಗಿ)}

\begin{sans}
pādāṅguṣṭhābhyāmūrdhvaṃ\endnote{ūrdhvaṃ~]  P, ūrdhva° M.} rajjuṃ dhṛtvā adhaḥ rajjuṃ hastābhyāṃ dhṛtvā ārohet\!। tṛṇajalūkaṃ bhavati\!॥101\,[102]\!॥
\end{sans}
%\bigskip

%102
\asana{Holding a [vertical] rope with the big toes above and the hands below, [the yogi] should climb up
it. [This] is the \engname{Caterpillar [Pose]} (\engiast{tṛṇajalūka}).}{ಇಳಿಬಿಟ್ಟಿರುವ ಹಗ್ಗವನ್ನು ಮೇಲ್ಭಾಗದಲ್ಲಿ ಕಾಲ್ಬೆರಳುಗಳಿಂದ ಮತ್ತು ಕೆಳಭಾಗದಲ್ಲಿ ಕೈಗಳಿಂದ ಹಿಡಿದು ಮೇಲಕ್ಕೆ ಏರಬೇಕು. ಇದು \tbf{ಕಂಬಳಿ ಹುಳುವಿನ ಭಂಗಿ}. (\tit{ತೃಣಜಲೂಕಾಸನ})}

%\newpage

\insimg[0.5]{HR_not_cropped/106_folio_28_recto_103.jpg}
\asananamestmp{vṛntāsanaṃ}{Grub Pose}{ವೃಂತಾಸನ (ಹುಳುವಿನ ಭಂಗಿ)}

\begin{sans}
ekayā muṣṭyā rajjuṃ dhṛtvā ārohet\!। vṛntāsanaṃ bhavati\!॥102\,[103]\!॥
\end{sans}
%\bigskip

%103
\asana{Having held a [vertical] rope with one fist, [the yogi] should climb up it. [This] is the \engname{Grub Pose}
(\engiast{vṛntāsana}).}{ಇಳಿಬಿಟ್ಟಿರುವ ಹಗ್ಗವನ್ನು ಕೇವಲ ಒಂದು ಮುಷ್ಟಿಯಿಂದ ಹಿಡಿದು ಮೇಲಕ್ಕೆ ಏರಬೇಕು. ಇದು \tbf{ಹುಳುವಿನ ಭಂಗಿ}. (\tit{ವೃಂತಾಸನ})}

\newpage

\insimg[0.56]{HR_not_cropped/107_folio_28_recto_104_E_ROPES.jpg}
\asananamestmp{krauñcāsanaṃ}{Crane Pose}{ಕ್ರೌಂಚಾಸನ (ಕೊಕ್ಕರೆಯ ಭಂಗಿ)}

\begin{sans}
ūrujānvantarābhyāṃ muṣṭī niṣkāsya tābhyāṃ rajjudvayaṃ dhṛtvā dantaiḥ bhāraṃ dhṛtvā ārohet\!। krauñcāsanaṃ bhavati\!॥103\,[104]\!॥
\end{sans}
%\bigskip

%104
\asana{Passing the fists through the thighs and knees, [the yogi] should hold two [vertical ropes] with
them, while holding a weight with the teeth. He should climb up. [This] is the \engname{Crane Pose}
(\engiast{krauñcāsana}).}{ತೊಡೆ ಮತ್ತು ಮೊಣಕಾಲುಗಳ ಸಂದಿನಿಂದ ಮುಷ್ಟಿಗಳನ್ನು ತೂರಿಸಿ ಇಳಿಬಿಟ್ಟಿರುವ ಎರಡು ಹಗ್ಗಗಳನ್ನು ಹಿಡಿದು, ಹಲ್ಲುಗಳಿಂದ ಒಂದು ಭಾರವನ್ನು ಕಚ್ಚಿ ಹಿಡಿದು ಮೇಲಕ್ಕೆ ಏರಬೇಕು. ಇದು \tbf{ಕೊಕ್ಕರೆಯ ಭಂಗಿ}. (\tit{ಕ್ರೌಂಚಾಸನ})}

\newpage

\insertFullImage{114_folio_30_recto_111.jpg}

\newpage

\pagesWithBorder

\null\vfill
{\leftskip=30pt
\begin{eng}
\prtheading[Poses that Pierce the Sun and Moon]{Poses that Pierce the Sun and Moon are [now] taught.}
\end{eng}
\begin{kan}
\prtheading[ಇನ್ನು ಮುಂದೆ ಸೂರ್ಯ ಮತ್ತು ಚಂದ್ರರನ್ನು ಭೇದಿಸುವ ಆಸನಗಳನ್ನು ಹೇಳಿಕೊಡಲಾಗುತ್ತದೆ]{ಇನ್ನು ಮುಂದೆ ಸೂರ್ಯ ಮತ್ತು ಚಂದ್ರರನ್ನು ಭೇದಿಸುವ ಆಸನಗಳನ್ನು ಹೇಳಿಕೊಡಲಾಗುತ್ತದೆ.}
\end{kan}\par}
\vfill

\newpage

\insimg[0.56]{HR_not_cropped/108_folio_28_verso_105_S_SUN_and_MOON.jpg}
\asananamestmp{varāhāsanaṃ}{Boar Pose}{ವರಾಹಾಸನ (ಹಂದಿಯ ಭಂಗಿ)}

\begin{sans}
sūryacandrabhedanānyāsanāni kathyante% Subsection heading (āsana group no. 6)

% new paragraph
kūrparau bhūmau sthāpayitvā jānubhyāmavanimavaṣṭabhya hastau mastake saṃsthāpya pārṣṇī\endnote{pārṣṇī~] \textit{corr.}~: pārṣṇi M.} nitambe saṃsthāpya tiṣṭhet\!। varāhāsanaṃ bhavati\!॥104\,[105]\!॥
\end{sans}
%\bigskip

%105
\begin{multicols}{2}
\begin{eng}
\Subsection*{Poses that Pierce the Sun and Moon are [now] taught.}
\end{eng}
\begin{kan}
\Subsection*{ಇನ್ನು ಮುಂದೆ ಸೂರ್ಯ ಮತ್ತು ಚಂದ್ರರನ್ನು ಭೇದಿಸುವ (ನಾಡಿಗಳ ಮೇಲೆ ಪ್ರಭಾವ ಬೀರುವ) ಆಸನಗಳನ್ನು ಹೇಳಿಕೊಡಲಾಗುತ್ತದೆ:}
\end{kan}
\end{multicols}

\asana{Having placed the elbows on the ground, [the yogi] should support [himself ] with the knees on
the ground, place the hands on the head and heels on the buttocks and remain thus. [This] is the
\engname{Boar Pose} (\engiast{varāhāsana}).}{ಮೊಣಕೈ ಮತ್ತು ಮೊಣಕಾಲುಗಳನ್ನು ನೆಲಕ್ಕೆ ಊರಿ, ತಲೆಯನ್ನು ಕೈಗಳ ಮೇಲಿಟ್ಟು, ಹಿಮ್ಮಡಿಗಳನ್ನು ಪೃಷ್ಠಕ್ಕೆ ತಾಗಿಸಿ ಹಾಗೆಯೇ ಇರಬೇಕು. ಇದು \tbf{ಹಂದಿಯ ಭಂಗಿ}. (\tit{ವರಾಹಾಸನ})}

\newpage

\insimg[0.5]{HR_not_cropped/109_folio_28_verso_106.jpg}
\asananamestmp{matsyendrapīṭhaṃ}{Matsyendra’s Seat}{ಮತ್ಸ್ಯೇಂದ್ರ ಪೀಠ}
%\vskip 1cm

\begin{sans}
vāmapārṣṇiṃ\endnote{°pārṣṇiṃ~] \textit{corr.}, °pārṣṇi M.} nābhau itarapādaṃ ūruṇi saṃsthāpya vāmahastena dakṣiṇajānuṃ\endnote{°jānuṃ~] P, °jānu M.} bahiḥpradeśena saṃveṣṭya vāmajānunaḥ adhaḥ pādāgraṃ dhṛtvā tiṣṭhet\!। matsyendrapīṭhaṃ\endnote{°pīṭhaṃ~] M(pc), pī+aṃ M(ac).} bhavati\!॥105\,[106]\!॥
\end{sans}
%\bigskip

%106
\asana{Placing the left heel on the navel [and] the other foot on the [opposite] thigh, [the yogi] should
clasp the outside of the right knee with the left hand and hold the toes [of the right foot, which
are] below the left knee. He should remain thus. [This] is \engname{Matsyendra’s Seat} (\engiast{matsyendrapīṭha}).}{ಎಡ ಹಿಮ್ಮಡಿಯನ್ನು ಹೊಕ್ಕುಳಿನ ಮೇಲಿಟ್ಟು, ಬಲ ಪಾದವನ್ನು ಎಡ ತೊಡೆಯ ಮೇಲಿರಿಸಬೇಕು. ಎಡಗೈಯಿಂದ ಬಲ ಮೊಣಕಾಲಿನ ಹೊರಭಾಗವನ್ನು ಅಪ್ಪಿ ಹಿಡಿದು, ಎಡ ಮೊಣಕಾಲಿನ ಕೆಳಗಿರುವ ಬಲಗಾಲಿನ ಬೆರಳುಗಳನ್ನು ಹಿಡಿಯಬೇಕು. ಇದು \tbf{ಮತ್ಸ್ಯೇಂದ್ರ ಪೀಠ}.}

%\newpage

\insimg[0.5]{HR_not_cropped/110_folio_29_recto_107.jpg}
\asananamestmp{yonyāsanaṃ}{Perineum Pose}{ಯೋನ್ಯಾಸನ}
%\vskip 1cm

\begin{sans}
pādau saṃmīlya pādāgre ādhāre pārṣṇī liṅgādadha ānīya pādatalayoḥ saṃviśet\!। yonyāsanaṃ bhavati\!॥106\,[107]\!॥
\end{sans}
%\bigskip

%107
\asana{[The yogi] should join the feet, bring the toes to the rectum and the heels below the penis, and sit
on the soles of the feet thus. [This] is the \engname{Perineum Pose} (\engiast{yonyāsana}).}{ಪಾದಗಳನ್ನು ಜೋಡಿಸಿ, ಕಾಲ್ಬೆರಳುಗಳನ್ನು ಗುದದ್ವಾರದ ಹತ್ತಿರ ಮತ್ತು ಹಿಮ್ಮಡಿಗಳನ್ನು ಶಿಶ್ನದ ಕೆಳಭಾಗಕ್ಕೆ ಬರುವಂತೆ ಇಟ್ಟು ಪಾದಗಳ ಮೇಲೆ ಕುಳಿತುಕೊಳ್ಳಬೇಕು. ಇದು \tbf{ಯೋನ್ಯಾಸನ}.}

\newpage

\insimg[0.54]{HR_not_cropped/111_folio_29_recto_108.jpg}
\asananamestmp{svastikāsanaṃ}{Lucky Mark Pose}{ಸ್ವಸ್ತಿಕಾಸನ}

\begin{sans}
ūruṇi pādatalaṃ saṃsthāpya itara ūruṇi itaraṃ niveśya saralaṃ tiṣṭhet\!। svastikāsanaṃ bhavati\!॥107\,[108]\!॥
\end{sans}
%\bigskip

%108
\asana{Placing the sole of one foot against the [opposite] thigh and fixing the other [sole] against the
other thigh, [the yogi] should sit upright. [This] is the \engname{Lucky Mark Pose} (\engiast{svastikāsana}).}{ಒಂದು ಪಾದದ ಅಡಿಯನ್ನು ಎದುರಿನ ತೊಡೆಗೆ ಮತ್ತು ಇನ್ನೊಂದು ಪಾದದ ಅಡಿಯನ್ನು ಮತ್ತೊಂದು ತೊಡೆಗೆ ತಾಗಿಸಿ ನೇರವಾಗಿ ಕುಳಿತುಕೊಳ್ಳಬೇಕು. ಇದು \tbf{ಸ್ವಸ್ತಿಕಾಸನ}.}

%\newpage

\insimg[0.54]{HR_not_cropped/112_folio_29_verso_109.jpg}
\asananamestmp{vajrāsanaṃ}{Thunderbolt Pose}{ವಜ್ರಾಸನ}

\begin{sans}
ekayā pārṣṇyā sīvanīṃ saṃpīḍya itarapārṣṇyā liṅgaṃ niṣpīḍya saralaṃ tiṣṭhet\!। vajrāsanaṃ bhavati\!॥108\,[109]\!॥
\end{sans}
%\bigskip

%109
\asana{Pressing the perineal seam with one heel and the penis with the heel of the other foot, [the yogi]
should sit upright. [This] is the \engname{Thunderbolt Pose} (\engiast{vajrāsana}).}{ಒಂದು ಹಿಮ್ಮಡಿಯಿಂದ ಶಿವನಿ ಭಾಗವನ್ನು ಮತ್ತು ಇನ್ನೊಂದು ಹಿಮ್ಮಡಿಯಿಂದ ಶಿಶ್ನದ ಮೇಲ್ಭಾಗವನ್ನು ಒತ್ತಿ ಹಿಡಿದು ನೇರವಾಗಿ ಕುಳಿತುಕೊಳ್ಳಬೇಕು. ಇದು \tbf{ವಜ್ರಾಸನ}.}

\newpage

\insimg[0.6]{HR_not_cropped/113_folio_29_verso_110.jpg}
\asananamestmp{utkaṭāsanaṃ}{Squatting Pose}{ಉತ್ಕಟಾಸನ}

\begin{sans}
pādatalābhyāṃ bhūmimavaṣṭabhya jānunī karṇamūle nayet\!। utkaṭāsanaṃ bhavati\!॥109\,[110]\!॥
\end{sans}
%\bigskip

%110
\asana{Supporting himself with the soles of the feet on the ground, [the yogi] should take his knees to the
base of his ears. [This] is the \engname{Squatting Pose} (\engiast{utkaṭāsana}).}{ಪಾದಗಳನ್ನು ನೆಲಕ್ಕೆ ಊರಿ ಕುಳಿತು, ಮೊಣಕಾಲುಗಳನ್ನು ಕಿವಿಗಳ ಬುಡಕ್ಕೆ ತಾಗಿಸಬೇಕು. ಇದು \tbf{ಉತ್ಕಟಾಸನ}.}

%\newpage

\insimg[0.6]{HR_not_cropped/114_folio_30_recto_111.jpg}
\asananamestmp{śuktyāsanaṃ}{Oyster Pose}{ಶುಕ್ತ್ಯಾಸನ (ಮುತ್ತಿನ ಚಿಪ್ಪಿನ ಭಂಗಿ)}

\begin{sans}
pārṣṇī\endnote{pārṣṇī~] \textit{corr.}, pārṣṇi M.} nābhau saṃsthāpya jaṅghā bahiḥpradeśena saṃmīlayet\!।  
śuktyāsanaṃ bhavati\!॥110\,[111]\!॥
\end{sans}
%\bigskip

%111
\asana{Placing the heels on the navel, [the yogi] should join the outsides of the lower legs. [This] is the \engname{Oyster Pose} (\engiast{śuktyāsana}).}{ಹಿಮ್ಮಡಿಗಳನ್ನು ಹೊಕ್ಕುಳಿನ ಹತ್ತಿರವಿಟ್ಟು, ಕಾಲುಗಳ ಕೆಳ ಭಾಗವನ್ನು ಪರಸ್ಪರ ಜೋಡಿಸಬೇಕು. ಇದು \tbf{ಮುತ್ತಿನ ಚಿಪ್ಪಿನ ಭಂಗಿ}. (\tit{ಶುಕ್ತ್ಯಾಸನ})}

\newpage

\insimg{HR_not_cropped/115_folio_30_recto_112.jpg}

\asananamestmp{śavāsanaṃ}{Corpse Pose}{ಶವಾಸನ}
\vskip 5mm

\begin{sans}
śavavadbhūmau tiṣṭhet\!। śavāsanaṃ bhavati\!॥111\,[112]\!॥
\end{sans}
\bigskip

%112
\asana{[The yogi] should remain like a corpse on the ground. [This] is the \engname{Corpse Pose} (\engiast{śavāsana}).}{ನೆಲದ ಮೇಲೆ ಶವದಂತೆ ನಿಶ್ಚಲವಾಗಿ ಮಲಗಬೇಕು. ಇದು \tbf{ಶವಾಸನ}.}

%\newpage
\bigskip

\insimg{HR_not_cropped/116_folio_30_verso_113.jpg}

\asananamestmp{uttānāsanaṃ}{Stretching Out [the Legs] Pose}{ಉತ್ತಾನಾಸನ}
\vskip 5mm

\begin{sans}
pādadvayaṃ vitanuyāt\!। uttānāsanaṃ\endnote{vitanuyāt~| uttānāsanaṃ~] \textit{conj.}~: vitanottānāsanaṃ M(pc), vitano+ttānāsanaṃ M(ac).} bhavati\!॥112\,[113]\!॥
\end{sans}
\bigskip

%113
\asana{[The yogi] should stretch out the legs. [This] is \engname{Stretching Out [the Legs] Pose} (\engiast{uttānāsana}).}{ಕಾಲುಗಳನ್ನು ಅಡ್ಡಲಾಗಿ ಪೂರ್ಣವಾಗಿ ಚಾಚಬೇಕು. ಇದು \tbf{ಉತ್ತಾನಾಸನ}.}

\newpage

\insimg[1]{HR_not_cropped/117_folio_30_verso_114_E_SUN_and_MOON.jpg}

\asananamestmp{sukhāsanaṃ}{Comfortable Pose}{ಸುಖಾಸನ}
\vskip 1cm

\begin{sans}
savyajānau\endnote{°jānau~] M(pc), °jānuṃ M(ac).} savyahastaṃ vāmajānau vāmahastaṃ ca sacinmudraṃ\endnote{sacin°~] \textit{corr.}, saciṃn° M.} nidhāya ardhonmīlitalocanastiṣṭhet\!। sukhāsanaṃ bhavati\!॥113\,[114]\!॥
\end{sans}
\bigskip

%114
\asana{Having placed in \textit{cinmudrā} the right hand on the right knee and the left hand on the left knee, [the
yogi] whose eyes are half open should remain thus. [This] is the \engname{Comfortable Pose} (\engiast{sukhāsana}).}{ಬಲಗೈಯನ್ನು ಬಲ ಮೊಣಕಾಲಿನ ಮೇಲೆ ಮತ್ತು ಎಡಗೈಯನ್ನು ಎಡ ಮೊಣಕಾಲಿನ ಮೇಲೆ ಚಿನ್ಮುದ್ರೆಯಲ್ಲಿರಿಸಿ, ಕಣ್ಣುಗಳನ್ನು ಅರ್ಧ ಮುಚ್ಚಿ ಕುಳಿತುಕೊಳ್ಳಬೇಕು. ಇದು \tbf{ಸುಖಾಸನ}}
